% concatenated from Diff_weighted_compositions_Bergman_spaces.tex
%% 
%% Copyright 2007-2025 Elsevier Ltd
%% 
%% This file is part of the 'Elsarticle Bundle'.
%% ---------------------------------------------
%% 
%% It may be distributed under the conditions of the LaTeX Project Public
%% License, either version 1.3 of this license or (at your option) any
%% later version.  The latest version of this license is in
%%    http://www.latex-project.org/lppl.txt
%% and version 1.3 or later is part of all distributions of LaTeX
%% version 1999/12/01 or later.
%% 
%% The list of all files belonging to the 'Elsarticle Bundle' is
%% given in the file `manifest.txt'.
%% 
%% Template article for Elsevier's document class `elsarticle'
%% with harvard style bibliographic references

\documentclass[preprint,12pt,times]{elsarticle}

%% Use the option review to obtain double line spacing
%\documentclass[preprint,review,12pt]{elsarticle}

%% Use the options 1p,twocolumn; 3p; 3p,twocolumn; 5p; or 5p,twocolumn
%% for a journal layout:
%\documentclass[final,1p,times]{elsarticle}
%% \documentclass[final,1p,times,twocolumn]{elsarticle}
 %\documentclass[final,3p,times]{elsarticle}
%% \documentclass[final,3p,times,twocolumn]{elsarticle}
 %\documentclass[final,5p,times]{elsarticle}
%% \documentclass[final,5p,times,twocolumn]{elsarticle}

%% For including figures, graphicx.sty has been loaded in
%% elsarticle.cls. If you prefer to use the old commands
%% please give \usepackage{epsfig}

%% The amssymb package provides various useful mathematical symbols
\usepackage{amssymb}
%% The amsmath package provides various useful equation environments.
\usepackage{amsmath}
%% The amsthm package provides extended theorem environments
\usepackage{amsthm}

%% The lineno packages adds line numbers. Start line numbering with
%% \begin{linenumbers}, end it with \end{linenumbers}. Or switch it on
%% for the whole article with \linenumbers.
%% \usepackage{lineno}



%\journal{}



\newtheorem{theorem}{Theorem}[section]
\newtheorem{corollary}[theorem]{Corollary}
\newtheorem{lemma}[theorem]{Lemma}
\newtheorem{proposition}[theorem]{Proposition}
\theoremstyle{definition}
\newtheorem{defn}[theorem]{Definition}
\theoremstyle{remark}
\newtheorem{remark}[theorem]{Remark}
\newtheorem*{example}{Example}
\numberwithin{equation}{section}
\usepackage[colorlinks, linkcolor=blue, anchorcolor=blue, citecolor=blue]{hyperref}
\usepackage{color}

\begin{document}

\begin{frontmatter}

%% Title, authors and addresses

%% use the tnoteref command within \title for footnotes;
%% use the tnotetext command for theassociated footnote;
%% use the fnref command within \author or \affiliation for footnotes;
%% use the fntext command for theassociated footnote;
%% use the corref command within \author for corresponding author footnotes;
%% use the cortext command for theassociated footnote;
%% use the ead command for the email address,
%% and the form \ead[url] for the home page:
%% \title{Title\tnoteref{label1}}
%% \tnotetext[label1]{}
%% \author{Name\corref{cor1}\fnref{label2}}
%% 
%% \ead[url]{home page}
%% \fntext[label2]{}
%% \cortext[cor1]{}
%% \affiliation{organization={},
%%             addressline={},
%%             city={},
%%             postcode={},
%%             state={},
%%             country={}}
%% \fntext[label3]{}

\title{Difference of weighted composition operators between some spaces of analytic functions} %% Article title

%% use optional labels to link authors explicitly to addresses:
%% \author[label1,label2]{}
%% \affiliation[label1]{organization={},
%%             addressline={},
%%             city={},
%%             postcode={},
%%             state={},
%%             country={}}
%%
%% \affiliation[label2]{organization={},
%%             addressline={},
%%             city={},
%%             postcode={},
%%             state={},
%%             country={}}
\author{Jiaoye Du\fnref{label1}}
\ead{djy18632952319@163.com}

\author{Cezhong Tong\fnref{label1}}
\ead{ctong@hebut.edu.cn; cezhongtong@hotmail.com}

%\author{Brett D. Wick\fnref{label2}}
%\ead{bwick@wustl.edu}

\author{Zicong Yang\fnref{label1}} %% Author name
\ead{zc25@hebut.edu.cn; zicongyang@126.com}

%\author{Zehua Zhou\fnref{label2}}
%\ead{zehuazhoumath@aliyun.com}
%% Author affiliation
%\cortext[cor1]{Corresponding author}
\affiliation[label1]{Department of Mathematics, Hebei University of Technology, Tianjin 300401, China}
%\affiliation[label2]{Department of Mathematics, Washington University in Saint Louis, Saint Louis, MO USA 63130}
%\affiliation[label2]{School of Mathematics, Tianjin University, Tianjin, 300354, China}

%% Abstract
\begin{abstract}
We first abtain a new and simpler proof of the main result in [IEOT, \textbf{93} (2021), 17], which characterized the bounded and compact differences of two weighted composition operators $C_{u,\varphi}-C_{v,\psi}$ acting between different Bergman spaces. More importantly, we get some characterizations for the difference of two weighted composition operators belonging to Schatten class. Futhermore, the compact difference of two weighted composition operators acting on Hardy-Hilbert spaces is also studied.
\end{abstract}

%%Graphical abstract
%\begin{graphicalabstract}
%%\includegraphics{grabs}
%\end{graphicalabstract}
%
%%%Research highlights
%\begin{highlights}
%\item Research highlight 1
%\item Research highlight 2
%\end{highlights}

%% Keywords
\begin{keyword}
%% keywords here, in the form: keyword \sep keyword
Bergman space, Weighted composition operator, Difference, Schatten class.
%% PACS codes here, in the form: \PACS code \sep code
\MSC[2020] 30H20; 47B33
%% MSC codes here, in the form: \MSC code \sep code
%% or \MSC[2008] code \sep code (2000 is the default)

\end{keyword}

\end{frontmatter}

%% Add \usepackage{lineno} before \begin{document} and uncomment 
%% following line to enable line numbers
%% \linenumbers

%% main text
%%
\section{Introduction}

Let $\mathbb{D}$ be the open unit disk in the complex plane $\mathbb{C}$ and $H(\mathbb{D})$ be the space of all analytic functions on $\mathbb{D}$. For $\alpha>-1$ and $0<p<\infty$, the standard weighted Bergman space $A_{\alpha}^p(\mathbb{D})$ consists of all functions $f\in H(\mathbb{D})$ such that
\[\|f\|_{A_{\alpha}^p}=\left(\int_{\mathbb{D}}|f(z)|^p{\rm d}A_{\alpha}(z)\right)^{1/p}\]
is finite. Here ${\rm d}A_{\alpha}(z)=(\alpha+1)(1-|z|^2)^{\alpha}{\rm d}A(z)$ and ${\rm d}A=\frac{1}{\pi}{\rm d}x{\rm d}y$ is the normalized Lebesgue measure on $\mathbb{D}$. It is known that $A_{\alpha}^p(\mathbb{D})$ is a Banach space for $1\leq p<\infty$. In particular, $A_{\alpha}^2(\mathbb{D})$ is a Hilbert space with the following inner product
\[\langle f, g\rangle=\int_{\mathbb{D}}f(z)\overline{g(z)}{\rm d}A_{\alpha}(z),\quad f,g\in A_{\alpha}^2(\mathbb{D}).\]
When $0<p<1$, the space $A_{\alpha}^p(\mathbb{D})$ is a complete metric space under the translation invariant metric $(f,g)\mapsto \|f-g\|_{A_{\alpha}^p}^p$.

The Hardy space $H^p(\mathbb{D})$ is defined by 
\[H^p(\mathbb{D})=\left\{f\in H(\mathbb{D}): \|f\|_{H^p}^p=\sup_{0<r<1}\int_{0}^{2\pi}|f(re^{i\theta})|^p\frac{{\rm d}\theta}{2\pi}<\infty\right\}.\]
Hardy space can be viewed as the limiting space of $A_{\alpha}^p(\mathbb{D})$ as $\alpha\to -1^+$. And the well-known Littlewood-Paley identity asserts that the $H^2$-norm can be converted to an area integral:
\begin{equation}\label{equa1.1}
\|f\|_{H^2}^2\simeq |f(0)|^2+\int_{\mathbb{D}}|f'(z)|^2(1-|z|^2){\rm d}A(z).
\end{equation}

The theory of analytic function spaces and operators acting on them, especially weighted composition operators, has witnessed an explosive growth in the past several decades. Denoted by $S(\mathbb{D})$ the set of all analytic self-maps of $\mathbb{D}$. Given $u\in H(\mathbb{D})$ and $\varphi\in S(\mathbb{D})$, the weighted composition operator $C_{u,\varphi}$ on $H(\mathbb{D})$ is defined by
\[C_{u,\varphi}f=u\cdot f\circ\varphi.\]
It is known that weighted composition operators are closely related to the isometries on Hardy and Bergman spaces, see for example \cite{13,14}. When $u=1$, it reduces to the composition operator $C_{\varphi}$. One can refer to the standard references \cite{9} and \cite{23} for various aspects on the theory of composition operators and refer to \cite{10,11} for the study of weighted composition operators on weighted Bergman spaces.

In the course of research on the topological structure of the space of all composition operators, which was initiated by Berkson \cite{2} and then by Shapiro and Sundberg \cite{24} in more details, problems of characterizing the path components and compact differences of two composition operators attracted broad interest. On Bergman spaces, the effort to characterize the compact difference of two composition operators was initiated by Moorhouse \cite{19} on $A_{\alpha}^2(\mathbb{D})$ and then by Saukko \cite{21, 22} from $A_{\alpha}^p(\mathbb{D})$ to $A_{\alpha}^q(\mathbb{D})$. By using Joint-Carleson measures, Koo and Wang \cite{15} studied the bounded and compact difference $C_{\varphi}-C_{\psi}$ on $A_{\alpha}^p$ over the unit ball. For more results about the difference of composition operators on various settings, see \cite{7,8,26} and references therein. 

Acharyya and Wu \cite{1} first considered compact differences of two weighted composition operators $C_{u,\varphi}-C_{v,\psi}$ between weighted Bergman spaces with the weights $u,v$ satisfying a certain growth condition. Recently, Choe et al. \cite{4} obtained complete characterizations for bounded and compact differences $C_{u,\varphi}-C_{v,\psi}$ on $A_{\alpha}^p(\mathbb{D})$ under an extra condition: $u,v\in A_{\alpha}^p(\mathbb{D})$. And then, they extended their results to operators acting from $A_{\alpha}^p(\mathbb{D})$ to $A_{\alpha}^q(\mathbb{D})$ under a condition that $u,v\in A_{\alpha}^q(\mathbb{D})$, see \cite{5}.

We first state the results in \cite{5}. To this end, we introduce several notations. For $\varphi,\psi\in S(\mathbb{D})$, put
\[\rho(z)=d(\varphi(z),\psi(z)),\quad z\in\mathbb{D},\]
where $d$ denotes the pseudo-hyperbolic distance on $\mathbb{D}$, see Section 2.1. Given a positive Borel measure $\mu$ on $\mathbb{D}$ and $\varphi\in S(\mathbb{D})$, the pull-back measure $\mu\circ\varphi^{-1}$ is defined by 
\[(\mu\circ\varphi)^{-1}(E)=\mu(\varphi^{-1}(E))\]
for any Borel set $E\subset\mathbb{D}$. For $u,v\in H(\mathbb{D})$, $\varphi,\psi\in S(\mathbb{D})$ and $\beta>0$, we now define several pull-back measures as follows:
\begin{align*}
\omega_{\varphi,u}^q:=\left(|\rho u|^q{\rm d}A_{\alpha}\right)\circ\varphi^{-1},\quad \omega_{\psi,v}^q:=\left(|\rho v|^q{\rm d}A_{\alpha}\right)\circ\psi^{-1},
\end{align*}
and
\begin{align*}
\sigma_{\varphi,\beta}^q:=\left[(1-\rho)^{\beta}|u-v|^q{\rm d}A_{\alpha}\right]\circ\varphi^{-1},\quad \sigma_{\psi,\beta}^q:=\left[(1-\rho)^{\beta}|u-v|^q{\rm d}A_{\alpha}\right]\circ\psi^{-1}.
\end{align*}
For simplicity, put
\[\omega_q=\omega_{\varphi,u}^q+\omega_{\psi,v}^q\quad {\rm and }\quad \sigma_{\beta,q}=\sigma_{\varphi,\beta}^q+\sigma_{\psi,\beta}^q.\]
The results in \cite{5} are described in terms of $(p,q)$-Bergman Carleson measures, see Section 2.3 for details about this notion.

\begin{theorem}[\cite{5}]\label{theorem1.1}
Let $\alpha>-1$, $0<p\leq q<\infty$ and $\beta>\frac{q}{p}(\alpha+2)$. Suppose $u,v\in A_{\alpha}^q(\mathbb{D})$ and $\varphi,\psi\in S(\mathbb{D})$. Then the following statements are equivalent.
\begin{itemize}
\item[(i)] $C_{u,\varphi}-C_{v,\psi}: A_{\alpha}^p(\mathbb{D})\to A_{\alpha}^q(\mathbb{D})$ is bounded (compact, resp.).
\item[(ii)] $\omega_q$ and $\sigma_{\beta,q}$ are (vanishing, resp.) $(p,q)$-Bergman Carleson measures.  
\end{itemize}
\end{theorem}

\begin{theorem}[\cite{5}]\label{theorem1.2}
Let $\alpha>-1$, $0<q<p<\infty$ and $\beta>\frac{q}{p}(\alpha+1+\max\{1,p\})$. Suppose $u,v\in A_{\alpha}^q(\mathbb{D})$ and $\varphi,\psi\in S(\mathbb{D})$. Then the following statements are equivalent.
\begin{itemize}
\item[(i)] $C_{u,\varphi}-C_{v,\psi}:A_{\alpha}^p(\mathbb{D})\to A_{\alpha}^q(\mathbb{D})$ is bounded;
\item[(ii)] $C_{u,\varphi}-C_{v,\psi}: A_{\alpha}^p(\mathbb{D})\to A_{\alpha}^q(\mathbb{D})$ is compact;
\item[(iii)] $\omega_q$ and $\sigma_{\beta,q}$ are $(p,q)$-Bergman Carleson measures.
\end{itemize}
\end{theorem}

In order to prove $(i)\Rightarrow(ii)$ in Theorem \ref{theorem1.1} and $(i)\Rightarrow(iii)$ in Theorem \ref{theorem1.2}, which are the hardest steps, several technical lemmas were applied in \cite{5}, see also \cite{4} for the case $p=q$. To the best of our knowledge, most of the studies of the difference of weighted composition operators rely on these lemmas in \cite{4} and \cite{5}, see for example \cite{16, 20}. Our first aim in this paper is to construct a new method and give a simpler proof for $(i)\Rightarrow (ii)$ in Theorem \ref{theorem1.1} and $(i)\Rightarrow(iii)$ in Theorem \ref{theorem1.2}. Moreover, the condition $u,v\in A_{\alpha}^q(\mathbb{D})$ is not needed in our proof. And the proof does not rely on those technical lemmas. Quite recently, Choe et al. \cite{3} extended the results in \cite{5} to the ball setting and removed the $L^q$-condition of the weights $u,v$. We point out that our method here is different from \cite{3}.

In addition to studying the bounded and compact differences of (weighted) composition operators, Choe   et al. \cite{6} characterized the Hilbert-Schmidt property of the difference $C_{\varphi}-C_{\psi} $ on $A_{\alpha}^2(\mathbb{D})$ and showed that $C_{\varphi}$ and $C_{\psi}$ are in the same path component in the space of composition operators under the Hilbert-Schmidt norm if and only if $C_{\varphi}-C_{\psi}$ is Hilbert-Schmidt. The result was then extended to the ball setting by Zhang and Zhou \cite{27}. Acharyya and Wu \cite{1} first characterized Hilbert-Schmidt difference of two weighted composition operators acting from Bergman space or Hardy space to an $L^2(\mu)$ space. For $0<p<\infty$, we recall that a compact operator $T$ acting on a separated Hilbert space $\mathcal{H}$ belongs to the Schatten class $S_p(\mathcal{H})$ if its sequence of singular numbers belongs to the sequence space $l^p$ (the singular numbers are the square roots of the eigenvalues of the positive operator $T^{*}T$). $T$ is said to be Hilbert-Schmidt if $T\in S_2(\mathcal{H})$. One can refer to \cite[Chapter 1]{29} for some details about Schatten classes. As far as we know, there are few papers studying the Schatten-$p$ class differences of weighted composition operators. So our second aim in this paper is to give some characterizations of operator differences belonging to Schatten-$p$ class.

For a positive Borel measure $\mu$ on $\mathbb{D}$ and $r>0$, we put
\[M_r(\mu)(z)=\frac{\mu(D(z,r))}{(1-|z|^2)^{2+\alpha}},\quad z\in\mathbb{D}\]
for the averaging function of $\mu$, where $D(z,r)$ is the Bergman metric disk centered at $z$ with radius $r$, see Section 2.1 and 2.4. Let 
\[{\rm d}\lambda(z)=\frac{{\rm d}A(z)}{(1-|z|^2)^2}\]
be the M$\ddot{\rm o}$bius invariant measure on $\mathbb{D}$. Our main result states as follows.

\begin{theorem}\label{theorem1.3}
Let $\alpha>-1$, $0<p<\infty$ and $\beta>0$. Suppose $u,v\in H(\mathbb{D})$ and $\varphi,\psi\in S(\mathbb{D})$. If $M_r(\omega_2)$ and $M_r(\sigma_{\beta,2})$ belong to $L^{\frac{p}{2}}(\mathbb{D},d\lambda)$ for some (or any) $r>0$, then $C_{u,\varphi}-C_{v,\psi}\in S_p(A_{\alpha}^2)$. Moreover, when $p\geq 2$ and $\beta\geq 2(\alpha+2)$, the above condition is also necessary for $C_{u,\varphi}-C_{v,\psi}\in S_p(A_{\alpha}^2)$.
\end{theorem}

By \eqref{equa1.1}, we know that $f\in H^2(\mathbb{D})$ if and only if $f'\in A_{1}^2(\mathbb{D})$. Choe et al. \cite{4} showed that $C_{\varphi}-C_{\psi}$ is compact on $H^2(\mathbb{D})$ if and only if $C_{\varphi',\varphi}-C_{\psi',\psi}$ is compact on $A_1^2(\mathbb{D})$ and obtained a compact criterion for the difference of two composition operators acting on $H^2(\mathbb{D})$. However, it will be more difficult to characterize the bounded and compact differences of two weighted composition operators $C_{u,\varphi}-C_{v,\psi}$ on Hardy spaces. Our next result shows a rigidity of the difference $C_{u,\varphi}-C_{v,\psi}$ on $H^2(\mathbb{D})$.
 
\begin{theorem}\label{theorem1.4}
Let $u,v\in H(\mathbb{D})$ and $\varphi,\psi\in S(\mathbb{D})$. Suppose $C_{u,\varphi}$ and $C_{v,\psi}$ are bounded on $H^2(\mathbb{D})$, then $C_{u,\varphi}-C_{v,\psi}$ is compact on $H^2(\mathbb{D})$ if and only if both $C_{u',\varphi}-C_{v',\psi}: H^2(\mathbb{D})\to A_1^2(\mathbb{D})$ and $C_{u\varphi',\varphi}-C_{v\psi',\psi}:A_1^2(\mathbb{D})\to A_1^2(\mathbb{D})$ are compact.
\end{theorem}

This paper is organized as follows. In Section 2, we collect some preliminary facts and auxiliary lemmas that will be used later. Section 3 is devoted to the new simpler proof of $(i)\Rightarrow(ii)$ in Theorem \ref{theorem1.1} and $(i)\Rightarrow(iii)$ in Theorem \ref{theorem1.2}. The proof of Theorem \ref{theorem1.3} is also presented in Section 4. Moreover, we obtain a Schatten-$p$ class characterization for the difference $C_{u,\varphi}-C_{v,\psi}$ via Reproducing Kernel Thesis. In Section 5, we prove Theorem \ref{theorem1.4} and characterize the linear sums of two composition operators on $H^2(\mathbb{D})$.

Throughout the paper we use the same letter $C$ to denote positive constants which may vary at different occurrences but do not depend on the essential argument. For non-negative quantities $A$ and $B$, we write $A\lesssim B$ (or equivalently $B\gtrsim A$) if there exists an absolute constant $C>0$ such that $A\leq CB$. $A\simeq B$ means both $A\lesssim B$ and $B\lesssim A$. 


\section{Preliminaries}

In this section, we recall some basic facts and present some auxiliary lemmas that will be used in the sequel.

\subsection{Pseudo-hyperbolic Distance}

The pseudo-hyperbolic distance between $z,w\in\mathbb{D}$ is given by 
\begin{equation*}
d(z,w)=\left|\frac{z-w}{1-\overline{z}w}\right|.
\end{equation*}
And the Bergman metric between $z,w\in\mathbb{D}$ is given by 
\begin{equation*}
\beta(z,w)=\frac{1}{2}\log\frac{1+d(z,w)}{1-d(z,w)}.
\end{equation*}
Write $D(z,r)=\{w\in\mathbb{D}: \beta(z,w)<r\}$ for the Bergman metric disk centered at $z$ with radius $r>0$.

It is easy to check that
\begin{equation}\label{equa2.1}
1-d^2(z,w)=\frac{(1-|z|^2)(1-|w|^2)}{|1-\overline{z}w|^2}.
\end{equation}

Given $r>0$, it is also well-known that
\begin{equation}\label{equa2.2}
A_{\alpha}(D(z,r))\simeq (1-|z|^2)^{2+\alpha},
\end{equation}
and 
\begin{equation}\label{equa2.3}
1-|z|^2\simeq 1-|w|^2\simeq |1-\overline{z}w|
\end{equation}
for all $z\in\mathbb{D}$ and $w\in D(z,r)$. Moreover,
\begin{equation}\label{equa2.4}
|1-\overline{a}z|\simeq |1-\overline{a}w|
\end{equation}
for all $a$, $z$ and $w$ in $\mathbb{D}$ with $\beta(z,w)<r$. Here, all the constants suppressed in these estimates depend only on $r$.


The following lemma, which is used in the proof of Theorem \ref{theorem1.3}, can be found in \cite[Lemma 2.2]{15}.

\begin{lemma}[\cite{15}]\label{lemma2.1}
Let $\alpha>-1$, $0<p<\infty$ and $0<r_1<r_2$. Then there is a constant $C=C(\alpha, p, r_1, r_2)$ such that
\begin{equation*}
|f(z)-f(w)|^p\leq C\frac{d^p(z,w)}{(1-|z|^2)^{2+\alpha}}\int_{D(z,r_2)}|f(\zeta)|^p{\rm d}A_{\alpha}(\zeta)
\end{equation*}
for any $f\in A_{\alpha}^p(\mathbb{D})$ and $z,w\in\mathbb{D}$ with $\beta(z,w)<r_1$.
\end{lemma}

A sequence $\{a_j\}$ in $\mathbb{D}$ is called an $r$-lattice in the Bergman metric if the following conditions hold:
\begin{itemize}
\item[(i)] $\bigcup\limits_{j=1}^{\infty}D(a_j,r)=\mathbb{D}$,
\item[(ii)] $\{D(a_j,r/2)\}_{j=1}^{\infty}$ are pairwise disjoint.
\end{itemize}
For any $r>0$, the existence of $r$-lattice in the Bergman metric is ensured by \cite[Lemma 4.8]{29}. And according to \cite[Lemma 4.7]{29}, with hypotheses (i) and (ii), it is easy to check that
\begin{itemize}
\item[(iii)] there exists a positive integer $N_0$ such that every point in $\mathbb{D}$ belongs to at most $N_0$ of the sets $D(a_j,4r)$. 
\end{itemize}

\subsection{Test Functions}

Recall the following pointwise estimate for derivatives of functions in $A_{\alpha}^p(\mathbb{D})$, see for example \cite[Lemma 2.1]{17}.

\begin{lemma}\label{lemma2.2}
Let $\alpha>-1$, $0<p<\infty$, $r>0$ and $n$ be an non-negative integer. There exists a constant $C=C(\alpha, p,r,n)>0$ such that
\begin{equation*}
|f^{(n)}(z)|^p\leq \frac{C}{(1-|z|^2)^{2+\alpha+np}}\int_{D(z,r)}|f(w)|^p{\rm d}A_{\alpha}(w)
\end{equation*}
for all $f\in A_{\alpha}^p(\mathbb{D})$ and $z\in\mathbb{D}$.
\end{lemma}

Note from Lemma \ref{lemma2.2} that each point evaluation of $n$th order: $f\mapsto f^{(n)}(z)$, is a continuous linear functional on $A_{\alpha}^p(\mathbb{D})$ for every $z\in\mathbb{D}$. Thus when $p=2$, by the Riesz representation theorem in Hilbert space theory, to each $z\in\mathbb{D}$ corresponds a unique function $K_z^{[n]}\in A_{\alpha}^2(\mathbb{D})$ such that
\begin{equation}\label{equa2.5}
f^{(n)}(z)=\langle f,K_z^{[n]}\rangle, \quad \forall f\in A_{\alpha}^2(\mathbb{D}).
\end{equation} 
$K_z^{[n]}$ is called the reproducing kernel function in $A_{\alpha}^2(\mathbb{D})$ at $z$ of order $n$, whose explicit formula is known as 
\begin{equation*}
K_z(w)=K_z^{[0]}(w)=\frac{1}{(1-\overline{z}w)^{2+\alpha}},\quad w\in\mathbb{D},
\end{equation*}
and $K_z^{[n]}(w)=\frac{\partial ^nK_z}{\partial\overline{z}^n}(w)$. It is known that
\begin{equation}\label{equa2.6}
\|K_z^{[n]}\|_{A_{\alpha}^2}^2=(K_z^{[n]})^{(n)}(z)\simeq \frac{1}{(1-|z|^2)^{2+\alpha+2n}}.
\end{equation}
Let $N,i$ be non-negative integers such that $N>\frac{2+\alpha}{p}$. For any $a\in\mathbb{D}$, let
\begin{equation*}
F_{a,N}^{[i]}(z)=\frac{z^i}{(1-\overline{a}z)^{N+i}},\quad z\in\mathbb{D}.
\end{equation*} 
According to \cite[Lemma 4.26]{29} and Forelli-Rudin's estimate, see \cite[Theorem 1.12]{28}, we have
\begin{equation*}
\|F_{a,N}^{[i]}\|_{A_{\alpha}^p}\simeq (1-|a|^2)^{\frac{2+\alpha}{p}-N-i}\quad {\rm and}\quad  \|F_{a,N}^{[i]}\|_{H^p}\simeq (1-|a|^2)^{\frac{1}{p}-N-i}.
\end{equation*}
Constants suppressed above are independent of $a$. We set 
\begin{equation*}
f_{a,N,p}^{[i]}(z)=(1-|a|^2)^{N+i-\frac{2+\alpha}{p}}F_{a,N}^{[i]}(z)
\end{equation*}
and
\begin{equation*}
g_{a,N,p}^{[i]}(z)=(1-|a|^2)^{N+i-\frac{1}{p}}F_{a,N}^{[i]}(z).
\end{equation*}
Then $\|f_{a,N,p}^{[i]}\|_{A_{\alpha}^p}\simeq 1$, $\|g_{a,N,p}^{[i]}\|_{H^p}\simeq 1$ and both $f_{a,N,p}^{[i]}$ and $g_{a,N,p}^{[i]}$ converges to 0 uniformly on compact subsets of $\mathbb{D}$ as $|a|\to 1^-$.

Besides, we have the following result, which derives from \cite[Theorem 4.33]{29}.

\begin{lemma}\label{lemma2.3}
Let $\{a_j\}$ be an $r$-lattice in the Bergman metric. Suppose $\alpha>-1$, $0<p<\infty$ and $N,i$ be non-negative integers with $N>\max\{1,\frac{1}{p}\}+\frac{1+\alpha}{p}$. For any $\mathbf{c}=\{c_j\}\in l^p$, let
\begin{equation*}
H_{\mathbf{c}, N,p}^{[i]}(z)=\sum_{j=1}^{\infty}c_j\frac{(1-|a_j|^2)^{N+i-\frac{2+\alpha}{p}}z^i}{(1-\overline{a}_jz)^{N+i}},\quad z\in\mathbb{D}.
\end{equation*}
Then $H_{\mathbf{c},N,p}^{[i]}\in A_{\alpha}^p(\mathbb{D})$ and $\|H_{\mathbf{c}, N,p}^{[i]}\|_{A_{\alpha}^p}\lesssim \|\{c_j\}\|_{l^p}$.
\end{lemma}

\subsection{Compact operator and Schatten-$p$ class}

Given two complete metrizable topological vector spaces $X$ and $Y$, a bounded linear operator $T: X\to Y$ is said to be compact if the image of each bounded sequence in $X$ has a convergent subsequence in $Y$.

The following compactness criterion derives from \cite[Lemma 2.1]{5}, see also \cite[Proposition 3.11]{9}.

\begin{lemma}\label{lemma2.4}
Let $\alpha\geq -1$ and $0<p,q<\infty$. Let $T$ be a linear combination of weighted composition operators and assume $T: A_{\alpha}^p(\mathbb{D})\to A_{\alpha}^q(\mathbb{D})$ is bounded. Then $T: A_{\alpha}^p(\mathbb{D})\to A_{\alpha}^q(\mathbb{D})$ is compact if and only if $\|Tf_n\|_{A_{\alpha}^q}\to 0$ for any bounded sequence in $A_{\alpha}^p(\mathbb{D})$ that converges to 0 uniformly on compact subsets of $\mathbb{D}$.
\end{lemma}

For any $z\in\mathbb{D}$, let $k_z^{[n]}=K_z^{[n]}/\|K_z^{[n]}\|_{A_{\alpha}^2}$. Suppose $T$ is a bounded linear operator on $A_{\alpha}^2$, we define a function 
\begin{equation*}
\widetilde{T}_{[n]}(z)=\langle Tk_z^{[n]}, k_z^{[n]}\rangle,\quad z\in\mathbb{D}.
\end{equation*}
$\widetilde{T}_{[n]}$ is called the $n$th-Berezin transform of $T$. When $n=0$, it reduces to the classcial Berezin transform.

\begin{lemma}\label{lemma2.5}
Suppose $T$ is a positive operator on $A_{\alpha}^2(\mathbb{D})$. If $T$ belongs to the trace class, i.e. $T\in S_1(A_{\alpha}^2)$, then $\widetilde{T}_{[n]}\in L^1(\mathbb{D},{\rm d}\lambda)$.
\end{lemma}

\begin{proof}
Fix an orthonormal basis $\{e_j\}$ for $A_{\alpha}^2(\mathbb{D})$. Since $T$ is positive, it belongs to the trace class if and only if
$$\sum_{j=1}^{\infty}\langle Te_j,e_j\rangle<\infty.$$
On the other hand, let $S=\sqrt{T}$. By \cite[Theorem 4.28]{29}, the reproducing formula \eqref{equa2.5} and \eqref{equa2.6}, we have
\begin{align*}
\|T\|_{S_1}&=\sum_{j=1}^{\infty}\langle Te_j,e_j\rangle=\sum_{j=1}^{\infty}\|Se_j\|_{A_{\alpha}^2}^2\\
&\gtrsim \sum_{j=1}^{\infty}\int_{\mathbb{D}}|(Se_j)^{(n)}(z)|^2(1-|z|^2)^{2n}{\rm d}A_{\alpha}(z)\\
&=\int_{\mathbb{D}}\left(\sum_{j=1}^{\infty}\langle Se_j, K_z^{[n]} \rangle^2\right)(1-|z|^2)^{2n}{\rm d}A_{\alpha}(z)\\
&=\int_{\mathbb{D}}\left(\sum_{j=1}^{\infty}\langle e_j, SK_z^{[n]}\rangle ^2\right)(1-|z|^2)^{2n}{\rm d}A_{\alpha}(z)\\
&=\int_{\mathbb{D}}\|SK_z^{[n]}\|_{A_{\alpha}^2}^2(1-|z|^2)^{2n}{\rm d}A_{\alpha}(z)\\
&=\int_{\mathbb{D}}\langle TK_z^{[n]}, K_z^{[n]}\rangle (1-|z|^2)^{2n}{\rm d}A_{\alpha}(z)\\
&\simeq \int_{\mathbb{D}}\langle Tk_z^{[n]}, k_z^{[n]}\rangle {\rm d}\lambda(z)=\int_{\mathbb{D}}\widetilde{T}_{[n]}(z){\rm d}{\lambda}(z).
\end{align*}
Therefore, $\int_{\mathbb{D}}\widetilde{T}_{[n]}(z){\rm d}\lambda(z)<\infty$. The proof is complete.
\end{proof}

\begin{lemma}\label{lemma2.6}
Suppose $T$ is a positive operator on $A_{\alpha}^2(\mathbb{D})$ and $1\leq p< \infty$. If $T\in S_p(A_{\alpha}^2)$, then $\widetilde{T}_{[n]}\in L^p(\mathbb{D},{\rm d}\lambda)$.
\end{lemma}

\begin{proof}
For $1\leq p<\infty$, by \cite[Proposition 1.31]{29}, we have
\[(\widetilde{T^p})_{[n]}(z)\geq \left(\widetilde{T}_{[n]}(z)\right)^p,\quad z\in\mathbb{D}.\]
Since $T$ is positive, $T\in S_p(A_{\alpha}^2)$ implies that $T^p$ is in the trace class. So it follows from Lemma \ref{lemma2.5} that
\begin{equation*}
\int_{\mathbb{D}}\left(\widetilde{T}_{[n]}(z)\right)^p{\rm d}\lambda(z)\leq\int_{\mathbb{D}}(\widetilde{T^p})_{[n]}(z){\rm d}\lambda(z)\lesssim \|T^p\|_{S_1}=\|T\|_{S_p}<\infty.
\end{equation*}
The proof is complete.
\end{proof}

\subsection{Carleson Measure}

Let $\mu$ be a positive Borel measure on $\mathbb{D}$. For $\alpha>-1$ and $0<p,q<\infty$, we say that $\mu$ is a $(p,q)$-Bergman Carleson measure if there exsits a constant $C>0$ such that
\begin{equation*}
\left(\int_{\mathbb{D}}|f(z)|^q{\rm d}\mu(z)\right)^{1/q}\leq C\|f\|_{A_{\alpha}^p}
\end{equation*}
for all $f\in A_{\alpha}^p(\mathbb{D})$. That is, $\mu$ is a $(p,q)$-Bergman Carleson measure if the embedding $A_{\alpha}^p(\mathbb{D})\subset L^q(\mathbb{D}, {\rm d}\mu)$ is continuous. If, in addition, the embedding $A_{\alpha}^p(\mathbb{D})\subset L^q(\mathbb{D},{\rm d}\mu)$ is compact, then $\mu$ is called a vanishing $(p,q)$-Bergman Carleson measure. Note that $(p,q)$-Bergman Carleson measures are finite.

For $r>0$ and $t>0$, we put
\[M_{r,t}(\mu)(z)=\frac{\mu(D(z,r))}{(1-|z|^2)^{t(2+\alpha)}},\quad z\in\mathbb{D}\]
for the weighted averaging function of $\mu$. When $t=1$, write $M_r(\mu):=M_{r,1}(\mu)$ for simplicity.

The geometric characterizations for $(p,q)$-Bergman Carleson measures have been known for some time, see \cite{18} and we state the results as follows.

\begin{lemma}\label{lemma2.7}
Let $0<p,q<\infty$ and $\mu$ be a positive Borel measure on $\mathbb{D}$.
\begin{itemize}
\item[(i)] If $0<p\leq q<\infty$, then $\mu$ is a $(p,q)$-Bergman Carleson measure if and only if 
$$\sup_{z\in\mathbb{D}}M_{r,\frac{q}{p}}(z)<\infty$$
for some (or any) $r>0$, and $\mu$ is a vanishing $(p,q)$-Bergman Carleson measure if and only if
$$\lim_{|z|\to 1^-}M_{r,\frac{q}{p}}(z)=0$$
for some (or any) $r>0$.
\item[(ii)] If $0<q<p<\infty$, then $\mu$ is a $(p,q)$-Bergman Carleson measure if and only if it is a vanishing $(p,q)$-Bergman Carleson measure if and only if
$$M_r(\mu)\in L^{\frac{p}{p-q}}(\mathbb{D},{\rm d}A_{\alpha})$$
for some (or any) $r>0$.
\end{itemize}
\end{lemma}

When $p=q$, we just call $\mu$ is a (vanishing) Bergman Carleson measure for simplicity, since the characterizations for $(p,p)$-Bergman Carleson measures are independent of $p$.

When $\mu$ is a Bergman Carleson measure and $r>0$, by Lemma \ref{lemma2.2} and Fubini's Theorem, it is easy to see that
\begin{equation}\label{equa2.7}
\int_{\mathbb{D}}|f(z)|^p{\rm d}\mu(z)\lesssim \int_{\mathbb{D}}|f(w)|^pM_r(\mu)(w){\rm d}A_{\alpha}(w)
\end{equation} 
for $0<p<\infty$ and any $f\in H(\mathbb{D})$. 

Given a positive Borel measure $\mu$ on $\mathbb{D}$, the Toeplitz operator $T_{\mu}$ acting on $A_{\alpha}^2(\mathbb{D})$ is the densely defined integral operator
\[T_{\mu}f(z)=\int_{\mathbb{D}}\frac{f(w)}{(1-z\overline{w})^{2+\alpha}}{\rm d}\mu(w),\quad z\in\mathbb{D}.\]
Clearly, $T_{\mu}$ is well defined on $H^{\infty}(\mathbb{D})$, the space of all bounded analytic functions on $\mathbb{D}$.

It is known that $T_{\mu}$ is bounded (compact, resp.) on $A_{\alpha}^2(\mathbb{D})$ if and only if $\mu$ is a (vanishing, resp.) Bergman Carleson measure. And for $0<p<\infty$, according to \cite[Theorem 7.18]{29}, $T_{\mu}$ is in $S_p(A_{\alpha}^2)$ if and only if 
\begin{equation}\label{equa2.8}
M_r(\mu)\in L^p(\mathbb{D},{\rm d}\lambda)
\end{equation}
for some (or any) $r>0$.

\begin{lemma}\label{lemma2.8}
Suppose $\mu$ is a vanishing Bergman Carleson measure and $R>r>0$. Let ${\rm d}\nu=M_r(\mu){\rm d}A_{\alpha}$. If $M_R(\mu)\in L^p(\mathbb{D},{\rm d}\lambda)$, then $T_\nu$ is in $S_p(A_{\alpha}^2)$.
\end{lemma}

\begin{proof}
By \eqref{equa2.3}, it is easy to see that $M_{R-r}(\nu)\lesssim M_R(\mu)$. So $M_{R-r}(\nu)\in L^{p}(\mathbb{D},{\rm d}\lambda)$, and then $T_{\nu}\in S_p(A_{\alpha}^2)$.
\end{proof}



\section{New proof of Theorems 1.1 and 1.2}

In this section, we present a new and simpler proof of $(i)\Rightarrow(ii)$ in Theorem \ref{theorem1.1} and $(i)\Rightarrow(iii)$ in Theorem \ref{theorem1.2}, which characterize the necessity for the boundedness and compactness of the difference $C_{u,\varphi}-C_{v,\psi}$ from $A_{\alpha}^p(\mathbb{D})$ to $A_{\alpha}^q(\mathbb{D})$.

\begin{proof}[{\bf New Proof of $(i)\Rightarrow(ii)$ in Theorem 1.1}]
Assume $p\leq q$ and $C_{u,\varphi}-C_{v,\psi}$ is bounded from $A_{\alpha}^p(\mathbb{D})$ to $A_{\alpha}^q(\mathbb{D})$. Then
\begin{equation*}
\sup_{a\in\mathbb{D}}\|(C_{u,\varphi}-C_{v,\psi})f_{a,N,p}^{[i]}\|_{A_{\alpha}^q}<\infty, \quad i=0,1.
\end{equation*}
Fix $r>0$. By \eqref{equa2.3}, we get
\begin{align*}
&\|(C_{u,\varphi}-C_{v,\psi})f_{a,N,p}^{[0]}\|_{A_{\alpha}^q}^q\\
&\quad \geq \int_{\varphi^{-1}(D(a,r))}\left|\frac{u(z)}{(1-\overline{a}\varphi(z))^N}-\frac{v(z)}{(1-\overline{a}\psi(z))^N}\right|^q(1-|a|^2)^{(N-\frac{2+\alpha}{p})q}{\rm d}A_{\alpha}(z)\\
&\quad \gtrsim\int_{\varphi^{-1}(D(a,r))}\left|u(z)-v(z)\left(\frac{1-\overline{a}\varphi(z)}{1-\overline{a}\psi(z)}\right)^N\right|^q\frac{1}{(1-|a|^2)^{\frac{q}{p}(2+\alpha)}}{\rm d}A_{\alpha}(z),
\end{align*}
and similarly,
\begin{align*}
&\|(C_{u,\varphi}-C_{v,\psi})f_{a,N,p}^{[1]}\|_{A_{\alpha}^q}^q\\
&\quad \gtrsim \frac{\int_{\varphi^{-1}(D(a,r))}\left|u(z)\varphi(z)-v(z)\psi(z)\left(\frac{1-\overline{a}\varphi(z)}{1-\overline{a}\psi(z)}\right)^{N+1}\right|^q{\rm d}A_{\alpha}(z)}{(1-|a|^2)^{\frac{q}{p}(\alpha+2)}}.
\end{align*}
Therefore,
\begin{equation}\label{equa3.1}
\sup_{a\in\mathbb{D}}\frac{\int_{\varphi^{-1}(D(a,r))}\left|u(z)-v(z)\left(\frac{1-\overline{a}\varphi(z)}{1-\overline{a}\psi(z)}\right)^N\right|^q{\rm d}A_{\alpha}(z)}{(1-|a|^2)^{\frac{q}{p}(\alpha+2)}}<\infty,
\end{equation}
and
\begin{equation}\label{equa3.2}
\sup_{a\in\mathbb{D}}\frac{\int_{\varphi^{-1}(D(a,r))}\left|u(z)\varphi(z)-v(z)\psi(z)\left(\frac{1-\overline{a}\varphi(z)}{1-\overline{a}\psi(z)}\right)^{N+1}\right|^q{\rm d}A_{\alpha}(z)}{(1-|a|^2)^{\frac{q}{p}(\alpha+2)}}<\infty.
\end{equation}
By \eqref{equa2.3}, $\left|\frac{1-\overline{a}\varphi(z)}{1-\overline{a}\psi(z)}\right|\lesssim 1$ when $\varphi(z)\in D(a,r)$. Multiplying the integrand in \eqref{equa3.1} by $|\psi(z)|^q\left|\frac{1-\overline{a}\varphi(z)}{1-\overline{a}\psi(z)}\right|^q$, adding it to \eqref{equa3.2}, and by the triangle inequality, we obtain
\begin{equation}\label{equa3.3}
\sup_{a\in\mathbb{D}}\frac{\int_{\varphi^{-1}(D(a,r))}|u(z)|^q\left|\frac{\varphi(z)-\psi(z)}{1-\overline{a}\psi(z)}\right|^q{\rm d}A_{\alpha}(z)}{(1-|a|^2)^{\frac{q}{p}(2+\alpha)}}<\infty.
\end{equation}
By \eqref{equa2.4}, $|1-\overline{a}\psi(z)|\simeq |1-\overline{\varphi(z)}\psi(z)|$ when $\varphi(z)\in D(a,r)$. So it follows from \eqref{equa3.3} that
\begin{align*}
\sup_{a\in\mathbb{D}}\frac{\int_{\varphi^{-1}(D(a,r))}|u(z)|^q\rho(z)^q{\rm d}A_{\alpha}(z)}{(1-|a|^2)^{\frac{q}{p}(\alpha+2)}}<\infty.
\end{align*}
This shows that $\omega_{\varphi,u}^q$ is a $(p,q)$-Bergman Carleson measure by Lemma \ref{lemma2.7}. Similarly, $\omega_{\psi,v}^q$ is also a $(p,q)$-Bergman Carleson measure.

On the other hand, when $\varphi(z)\in D(a,r)$, we have
\begin{equation}\label{equa3.4}
\begin{split}
&|u(z)-v(z)|\left|\frac{1-\overline{a}\varphi(z)}{1-\overline{a}\psi(z)}\right|^N\\
&\quad \leq \left|u(z)-v(z)\left(\frac{1-\overline{a}\varphi(z)}{1-\overline{a}\psi(z)}\right)^N\right|+|u(z)|\left|1-\left(\frac{1-\overline{a}\varphi(z)}{1-\overline{a}\psi(z)}\right)^N\right|\\
&\quad \lesssim \left|u(z)-v(z)\left(\frac{1-\overline{a}\varphi(z)}{1-\overline{a}\psi(z)}\right)^N\right|+|u(z)|\left|\frac{\varphi(z)-\psi(z)}{1-\overline{a}\psi(z)}\right|
\end{split}
\end{equation}
And by \eqref{equa2.1}, \eqref{equa2.3} and \eqref{equa2.4}, we have
\begin{equation}\label{equa3.5}
\left|\frac{1-\overline{a}\varphi(z)}{1-\overline{a}\psi(z)}\right|\gtrsim \frac{(1-|\varphi(z)|^2)(1-|\psi(z)|^2)}{|1-\overline{\varphi(z)}\psi(z)|^2}=1-\rho^2(z)
\end{equation}
for $z\in \varphi^{-1}(D(a,r))$. So we combine \eqref{equa3.1} and \eqref{equa3.3} to obtain
\begin{equation*}
\sup_{a\in\mathbb{D}}\frac{\int_{\varphi^{-1}(D(a,r))}|u(z)-v(z)|^q(1-\rho(z))^{qN}{\rm d}A_{\alpha}(z)}{(1-|a|^2)^{\frac{q}{p}(2+\alpha)}}<\infty.
\end{equation*}
This means that $\sigma_{\varphi,\beta}^q$ is a $(p,q)$-Bergman Carleson measure for $\beta\geq qN>\frac{q}{p}(2+\alpha)$ by Lemma \ref{lemma2.7}. Similarly, $\sigma_{\psi,\beta}^q$ is also a $(p,q)$-Bergman Carleson measure.

Finally, the proof for the compactness part is just a modification. By Lemma \ref{lemma2.4}, one only needs to write ``$\lim\limits_{|a|\to 1}$'' instead of ``$\sup\limits_{a\in\mathbb{D}}$'' and write ``$\to 0$'' instead of ``$<\infty$'', respectively. We omit the routine details.
\end{proof}


According to the above proof, we conclude the following corollary, which coincides with \cite[Theorem 1.1]{3}.

\begin{corollary}
Let $\alpha>-1$, $0<p\leq q<\infty$ and $N>\frac{2+\alpha}{p}$. Suppose $u,v\in H(\mathbb{D})$ and $\varphi,\psi\in S(\mathbb{D})$. Then
\begin{itemize}
\item[(i)] $C_{u,\varphi}-C_{v,\psi}: A_{\alpha}^p(\mathbb{D})\to A_{\alpha}^q(\mathbb{D})$ is bounded if and only if 
\begin{equation*}
\sup_{a\in\mathbb{D}}\|(C_{u,\varphi}-C_{v,\psi})f_{a,N,p}^{[i]}\|_{A_{\alpha}^q}<\infty, \quad i\in \mathbb{N}.
\end{equation*}
\item[(ii)] 
$C_{u,\varphi}-C_{v,\psi}: A_{\alpha}^p(\mathbb{D})\to A_{\alpha}^q(\mathbb{D})$ is compact if and only if 
\begin{equation*}
\lim_{|a|\to 1}\|(C_{u,\varphi}-C_{v,\psi})f_{a,N,p}^{[i]}\|_{A_{\alpha}^q}=0, \quad i\in \mathbb{N}.
\end{equation*}
\end{itemize}
\end{corollary}

Now we proceed to the proof of $(i)\Rightarrow (iii)$ in Theorem \ref{theorem1.2}. Before that, we recall some basic facts about the Khinchine's inequality, which is an important tool in complex and functional analysis. An introduction to this topic can be found in Appendix of \cite{12}.

Let $\{r_j(t)\}$ denotes the sequence of Rademacher functions defined by
$$r_0(t)=\begin{cases}1, &\mbox{if}\ 0\leq t-[t]<\frac{1}{2}\\-1, &\mbox{if}\ \frac{1}{2}\leq t-[t]<1,\end{cases}$$ 
where $[t]$ denotes the largest integer not greater than $t$ and $r_j(t)=r_0(2^jt)$ for $j=1, 2,\cdots$.
If $0<p<\infty$, then Khinchine's inequality states that
$$\left(\sum_{j}\left|c_{j}\right|^{2}\right)^{p / 2} \simeq \int_{0}^{1}\left|\sum_{j} c_{j} r_{j}(t)\right|^{p} {\rm d}t$$
for complex sequences $\left\{c_{j}\right\}$

\begin{proof}[{\bf New proof of $(i)\Rightarrow (iii)$ in Theorem 1.2}]
Assume $0<q<p<\infty$ and $C_{u,\varphi}-C_{v,\psi}: A_{\alpha}^p(\mathbb{D})\to A_{\alpha}^q(\mathbb{D})$ is bounded. Fix $r>0$. Let $\{a_j\}$ be an $r$-lattice in the Bergman metric and $\mathbf{c}=\{c_j\}\in l^p$. For $N>\max\{1,\frac{1}{p}\}+\frac{1+\alpha}{p}$, we use the test function $H_{\mathbf{c},N,p}^{[i]}$ defined in Lemma \ref{lemma2.3} to obtain
\begin{equation*}
\|(C_{u,\varphi}-C_{v,\psi})H_{\mathbf{c},N,p}^{[i]}\|_{A_{\alpha}^q}\lesssim \|\{c_j\}\|_{l^p}, \quad i=0,1.
\end{equation*}
For $i=0$, we obtain
\begin{equation}\label{equa3.6}
\begin{split}
&\int_{\mathbb{D}}\left|\sum_{j=1}^{\infty}c_j\left(u(z)\frac{(1-|a_j|^2)^{N-\frac{2+\alpha}{p}}}{(1-\overline{a}_j\varphi(z))^N}-v(z)\frac{(1-|a_j|^2)^{N-\frac{2+\alpha}{p}}}{(1-\overline{a}_j\psi(z))^N}\right)\right|^q{\rm d}A_{\alpha}(z)\\
&\quad \lesssim \|\{c_j\}\|_{l^p}^q.
\end{split}
\end{equation}
In \eqref{equa3.6}, we replace $c_j$ by $r_j(t)c_j$, so that the right hand side does not change. Then intergrating both sides with respect to $t$ from 0 to 1, and by Khinchine's inequaltiy and Fubini's Theorem, we obtain
\begin{equation*}
\begin{split}
&\int_{\mathbb{D}}\left(\sum_{j=1}^{\infty}|c_j|^2\left|u(z)\frac{(1-|a_j|^2)^{N-\frac{2+\alpha}{p}}}{(1-\overline{a}_j\varphi(z))^N}-v(z)\frac{(1-|a_j|^2)^{N-\frac{2+\alpha}{p}}}{(1-\overline{a}_j\psi(z))^N}\right|^2\right)^{\frac{q}{2}}{\rm d}A_{\alpha}(z)\\
&\simeq\int_{\mathbb{D}}\int_0^1\left|\sum_{j=1}^{\infty}c_jr_j(t)\left(u(z)\frac{(1-|a_j|^2)^{N-\frac{2+\alpha}{p}}}{(1-\overline{a}_j\varphi(z))^N}-v(z)\frac{(1-|a_j|^2)^{N-\frac{2+\alpha}{p}}}{(1-\overline{a}_j\psi(z))^N}\right)\right|^q{\rm d}t{\rm d}A_{\alpha}(z)\\
&=\int_{0}^1\int_{\mathbb{D}}\left|\sum_{j=1}^{\infty}c_jr_j(t)\left(u(z)\frac{(1-|a_j|^2)^{N-\frac{2+\alpha}{p}}}{(1-\overline{a}_j\varphi(z))^N}-v(z)\frac{(1-|a_j|^2)^{N-\frac{2+\alpha}{p}}}{(1-\overline{a}_j\psi(z))^N}\right)\right|^q{\rm d}A_{\alpha}(z){\rm d}t\\
&\lesssim \|\{c_j\}\|_{l^p}^q.
\end{split}
\end{equation*}
By \eqref{equa2.3}, $\chi_{D(a_j,2r)}(\varphi(z))\lesssim \frac{1-|a_j|^2}{|1-\overline{a}_j\varphi(z)|}$. It follows that
\begin{equation*}
\begin{split}
\int_{\mathbb{D}}\left(\sum_{j=1}^{\infty}\frac{|c_j|^2\chi_{D(a_j,2r)}(\varphi(z))\left|u(z)-v(z)\left(\frac{1-\overline{a}_j\varphi(z)}{1-\overline{a}_j\psi(z)}\right)^N\right|^2}{(1-|a_j|^2)^{2(2+\alpha)/p}}\right)^{q/2}{\rm d}A_{\alpha(z)} \lesssim \|\{c_j\}\|_{l^p}^q.
\end{split}
\end{equation*}
Recall that there is a positive integer $N_0$ such that each point $z\in\mathbb{D}$ belongs to at most $N_0$ of the disks $D(a_j,2r)$. Applying Minkowski's inequality if $\frac{2}{q}\leq 1$ and H\"older's inequality if $\frac{2}{q}>1$, we obtain
\begin{equation*}
\begin{split}
&\sum_{j=1}^{\infty}|c_j|^q\frac{\int_{\varphi^{-1}(D(a_j,2r))}\left|u(z)-v(z)\left(\frac{1-\overline{a}_j\varphi(z)}{1-\overline{a}_j\psi(z)}\right)^N\right|^q{\rm d}A_{\alpha}(z)}{(1-|a_j|^2)^{q(2+\alpha)/p}}\\
&\quad\leq C\int_{\mathbb{D}}\left(\sum_{j=1}^{\infty}\frac{|c_j|^2\chi_{D(a_j,2r)}(\varphi(z))\left|u(z)-v(z)\left(\frac{1-\overline{a}\varphi(z)}{1-\overline{a}\psi(z)}\right)^N\right|^2}{(1-|a_j|^2)^{2(2+\alpha)/p}}\right)^{q/2}{\rm d}A_{\alpha(z)}\\
&\quad \lesssim \|\{c_j\}\|_{l^p}^q,
\end{split}
\end{equation*}
where $C=\max\{N_0^{1-\frac{q}{2}},1\}$. This implies that the sequence 
\[\left\{\frac{\int_{\varphi^{-1}(D(a_j,2r))}\left|u(z)-v(z)\left(\frac{1-\overline{a}_j\varphi(z)}{1-\overline{a}_j\psi(z)}\right)^N\right|^q{\rm d}A_{\alpha}(z)}{(1-|a_j|^2)^{q(2+\alpha)/p}}\right\}_{j=1}^{\infty}\]
belongs to the dual of $l^{\frac{p}{q}}$ or, equivalently,
\begin{equation}\label{equa3.7}
\sum_{j=1}^{\infty}\left(\frac{\int_{\varphi^{-1}(D(a_j,2r))}\left|u(z)-v(z)\left(\frac{1-\overline{a}_j\varphi(z)}{1-\overline{a}_j\psi(z)}\right)^N\right|^q{\rm d}A_{\alpha}(z)}{(1-|a_j|^2)^{q(2+\alpha)/p}}\right)^{\frac{p}{p-q}}<\infty.
\end{equation}
Similarly, taking $i=1$, we obtain
\begin{equation}\label{equa3.8}
\sum_{j=1}^{\infty}\left(\frac{\int_{\varphi^{-1}(D(a_j,2r))}\left|u(z)\varphi(z)-v(z)\psi(z)\left(\frac{1-\overline{a}_j\varphi(z)}{1-\overline{a}_j\psi(z)}\right)^{N+1}\right|^qdA_{\alpha}(z)}{(1-|a_j|^2)^{q(2+\alpha)/p}}\right)^{\frac{p}{p-q}}<\infty.
\end{equation}
By \eqref{equa2.3}, $\left|\frac{1-\overline{a}_j\varphi(z)}{1-\overline{a}_j\psi(z)}\right|\lesssim 1$ when $\varphi(z)\in D(a,r)$ and constant suppressed is independent of $a_j$, multiplying the integrand in \eqref{equa3.7} by $|\psi(z)|^q\left|\frac{1-\overline{a_j}\varphi(z)}{1-\overline{a}_j\psi(z)}\right|^q$, adding it to \eqref{equa3.8} and by the triangle inequality and \eqref{equa2.4}, we get
\begin{equation}\label{equa3.9}
\sum_{j=1}^{\infty}\left(\frac{\int_{\varphi^{-1}(D(a_j,2r))}|u(z)|^q\rho(z)^q{\rm d}A_{\alpha}(z)}{(1-|a_j|^2)^{q(2+\alpha)/p}}\right)^{\frac{p}{p-q}}<\infty.
\end{equation}
Therefore, by \eqref{equa2.2} and \eqref{equa2.3}, we obtain
\begin{equation}\label{equa3.10}
\begin{split}
&\int_{\mathbb{D}}\left(\frac{\int_{\varphi^{-1}(D(w,r))}|u(z)|^q\rho(z)^q{\rm d}A_{\alpha}(z)}{(1-|w|^2)^{2+\alpha}}\right)^{\frac{p}{p-q}}{\rm d}A_{\alpha}(w)\\
&\quad \leq \sum_{j=1}^{\infty}\int_{D(a_j,r)}\left(\frac{\int_{\varphi^{-1}(D(w,r))}|u(z)|^q\rho(z)^q{\rm d}A_{\alpha}(z)}{(1-|w|^2)^{2+\alpha}}\right)^{\frac{p}{p-q}}{\rm d}A_{\alpha}(w)\\
&\quad \lesssim \sum_{j=1}^{\infty}\left(\frac{\int_{\varphi^{-1}(D(a_j,2r))}|u(z)|^q\rho(z)^q{\rm d}A_{\alpha}(z)}{(1-|a_j|^2)^{2+\alpha}}\right)^{\frac{p}{p-q}}A_{\alpha}(D(a_j,2r))\\
&\quad \lesssim \sum_{j=1}^{\infty}\left(\frac{\int_{\varphi^{-1}(D(a_j,2r))}|u(z)|^q\rho(z)^q{\rm d}A_{\alpha}(z)}{(1-|a_j|^2)^{q(2+\alpha)/p}}\right)^{\frac{p}{p-q}}<\infty.
\end{split}
\end{equation}
This implies that $\omega_{\varphi,u}^q$ is a $(p,q)$-Bergman Carleson measure by Lemma \ref{lemma2.7}. Similarly, $\omega_{\psi,v}^q$ is also a $(p,q)$-Bergman Carleson measure.

On the other hand, when $\varphi(z)\in D(a_j,2r)$, by \eqref{equa3.4}, \eqref{equa3.5} and \eqref{equa2.4}, we have
\begin{equation*}
\begin{split}
|u(z)-v(z)|(1-\rho(z))^N&\lesssim |u(z)-v(z)|\left|\frac{1-\overline{a}_j\varphi(z)}{1-\overline{a}_j\psi(z)}\right|^N\\
&\lesssim \left|u(z)-v(z)\left(\frac{1-\overline{a}_j\varphi(z)}{1-\overline{a}_j\psi(z)}\right)^N\right|+|u(z)|\rho(z).
\end{split}
\end{equation*}
So we combine \eqref{equa3.7} and \eqref{equa3.9} to obtain
\begin{equation*}
\sum_{j=1}^{\infty}\left(\frac{\int_{\varphi^{-1}(D(a_j,2r))}|u(z)-v(z)|^q(1-\rho(z))^{qN}{\rm d}A_{\alpha}(z)}{(1-|a_j|^2)^{q(2+\alpha)/p}}\right)^{\frac{p}{p-q}}<\infty.
\end{equation*}
Then through a similar argument as \eqref{equa3.10}, we get
\begin{equation*}
\int_{\mathbb{D}}\left(\frac{\int_{\varphi^{-1}(D(w,r))}|u(z)-v(z)|^q(1-\rho(z))^{qN}{\rm d}A_{\alpha}(z)}{(1-|w|^2)^{2+\alpha}}\right)^{\frac{p}{p-q}}{\rm d}A_{\alpha}(w)<\infty.
\end{equation*}
This implies that $\sigma_{\varphi,\beta}^q$ is a $(p,q)$-Bergman Carleson measure for $\beta\geq qN>\frac{q}{p}(\alpha+1+\max\{1,p\})$ by Lemma \ref{lemma2.7}. Similarly, $\sigma_{\psi,\beta}^q$ is also a $(p,q)$-Bergman Carleson measure. The proof is now complete.
\end{proof}

According to the above proof, we conclude the following corollary.

\begin{corollary}\label{corollary3.2}
Let $\alpha>-1$, $0<q<p<\infty$ and $N>\max\{1,\frac{1}{p}\}+\frac{1+\alpha}{p}$. Suppose $u,v\in H(\mathbb{D})$ and $\varphi,\psi\in S(\mathbb{D})$. Then $C_{u,\varphi}-C_{v,\psi}: A_{\alpha}^p(\mathbb{D})\to A_{\alpha}^q(\mathbb{D})$ is bounded if and only it is compact, if and only if
\[\|(C_{u,\varphi}-C_{v,\psi})H_{{\mathbf{c}},N,p}^{[i]}\|_{A_{\alpha}^q}\lesssim \|\{c_j\}\|_{l^p},\quad \forall i\in\mathbb{N}\]
for any $\mathbf{c}=\{c_j\}\in l^p$.
\end{corollary}


\section{Proof of Theorem 1.3}

In this section, we aim to prove Theorem \ref{theorem1.3}, which characterizes the Schatten class difference of weighted composition operators on $A_{\alpha}^2(\mathbb{D})$. And then we obtain some direct corollaries.

\begin{proof}[{\bf Proof of Theorem 1.3}]
{\it Sufficiency}: Fix $r>0$, set 
$$G_r=\{z\in\mathbb{D}:\beta(\varphi(z),\psi(z))<r\}.$$
For any $f\in A_{\alpha}^2(\mathbb{D})$, we have
\begin{equation*}
\begin{split}
\langle (C_{u,\varphi}-C_{v,\psi})^*(C_{u,\varphi}-C_{v,\psi})f,f\rangle&=\|(C_{u,\varphi}-C_{v,\psi})f\|_{A_{\alpha}^2}^2\\
&=\int_{\mathbb{D}\backslash G_r}+\int_{G_r}|u\cdot f\circ\varphi-v\cdot f\circ\psi|^2{\rm d}A_{\alpha}\\
&:=I(f)+II(f).
\end{split}
\end{equation*}
Now we estimate each integral separately. For the integral $I(f)$, note that $\rho(z)\geq \tanh(r)$ for $z\in \mathbb{D}\backslash G_r$. So by \eqref{equa2.7}, we obtain
\begin{equation}\label{equa4.1}
\begin{split}
I(f)&\lesssim \int_{\mathbb{D}\backslash G_r}\left(|u\cdot f\circ\varphi|^2+|v\cdot f\circ\psi|^2\right){\rm d}A_{\alpha}\\
&\lesssim \int_{\mathbb{D}}\left(|u\rho|^2|f\circ\varphi|^2+|v\rho|^2|f\circ\psi|^2\right){\rm d}A_{\alpha}\\
&\lesssim \int_{\mathbb{D}}|f|^2M_{2r}(\omega_2){\rm d}A_{\alpha}.
\end{split}
\end{equation}
Meanwhile, clearly,
\begin{align*}
II(f)&\lesssim \int_{G_r}|u-v|^2(|f\circ\varphi|^2+|f\circ\psi|^2)+(|u|^2+|v|^2)|f\circ\varphi-f\circ\psi|^2{\rm d}A_{\alpha}\\
&:=II_1(f)+II_2(f).
\end{align*}
For the integral $II_1(f)$, note that $1-\rho(z)\geq 1-\tanh(r)$ for $z\in G_r$, so by \eqref{equa2.7}, we obtain
\begin{equation}\label{equa4.2}
\begin{split}
II_1(f)&\lesssim \int_{G_r}|u-v|^2(1-\rho)^{\beta}(|f\circ\varphi|^2+|f\circ\psi|^2){\rm d}A_{\alpha}\\
&\lesssim \int_{\mathbb{D}}|f|^2M_{2r}(\sigma_{\beta,2}){\rm d}A_{\alpha}.
\end{split}
\end{equation}
For the integral $II_2(f)$, applying Lemma \ref{lemma2.1}, when $z\in G_r$, we have
\begin{equation*}
|f(\varphi(z))-f(\psi(z))|^2\lesssim \frac{\rho(z)^2}{(1-|\varphi(z)|^2)^{2+\alpha}}\int_{D(\varphi(z),2r)}|f(w)|^2{\rm d}A_{\alpha}(w),
\end{equation*}
and
\begin{equation*}
|f(\varphi(z))-f(\psi(z))|^2\lesssim \frac{\rho(z)^2}{(1-|\psi(z)|^2)^{2+\alpha}}\int_{D(\psi(z),2r)}|f(w)|^2{\rm d}A_{\alpha}(w).
\end{equation*}
Then we use Fubini's Theorem to obtain
\begin{equation}\label{equa4.3}
\begin{split}
II_2(f)&\lesssim\int_{G_r}\frac{|u(z)|^2\rho(z)^2}{(1-|\varphi(z)|^2)^{2+\alpha}}\left(\int_{D(\varphi(z),2r)}|f(w)|^2{\rm d}A_{\alpha}(w)\right){\rm d}A_{\alpha}(z)\\
&\quad +\int_{G_r}\frac{|v(z)|^2\rho(z)^2}{(1-|\psi(z)|^2)^{2+\alpha}}\left(\int_{D(\psi(z),2r)}|f(w)|^2{\rm d}A_{\alpha}(w)\right){\rm d}A_{\alpha}(z)\\
&\lesssim \int_{\mathbb{D}}|f|^2M_{2r}(\omega_2){\rm d}A_{\alpha}.
\end{split}
\end{equation}

Let ${\rm d}\mu=(M_{2r}(\omega_2)+M_{2r}(\sigma_{\beta,2})){\rm d}A_{\alpha}$. Combining \eqref{equa4.1}, \eqref{equa4.2} and \eqref{equa4.3}, we obtain
\[(C_{u,\varphi}-C_{v,\psi})^*(C_{u,\varphi}-C_{v,\psi})\lesssim T_{\mu}.\]
By Lemma \ref{lemma2.8}, the assumption implies that $T_{\mu}\in S_{\frac{p}{2}}(A_{\alpha}^2)$. It follows that $(C_{u,\varphi}-C_{v,\psi})^*(C_{u,\varphi}-C_{v,\psi})\in S_{\frac{p}{2}}(A_{\alpha}^2)$, and then $C_{u,\varphi-}C_{v,\psi}\in S_p(A_{\alpha}^2)$.

{\it Necessity}: Assume $C_{u,\varphi}-C_{v,\psi}\in S_p(A_{\alpha}^2)$ and $p\geq 2$. Then $(C_{u,\varphi}-C_{v,\psi})^*(C_{u,\varphi}-C_{v,\psi})\in S_{\frac{p}{2}}(A_{\alpha}^2)$. And it follows from Lemma \ref{lemma2.6} that the function
\[z\mapsto \|(C_{u,\varphi}-C_{v,\psi})k_z^{[n]}\|_{A_{\alpha}^2}^2\]
belongs to $L^{\frac{p}{2}}(\mathbb{D},{\rm d}\lambda)$.

For any $r>0$, by \eqref{equa2.3} and \eqref{equa2.6}, we have
\begin{equation}\label{equa4.4}
\begin{split}
&\|(C_{u,\varphi}-C_{v,\psi})k_z^{[0]}\|_{A_{\alpha}^2}^2\\
&\gtrsim\int_{\varphi^{-1}(D(z,r))}\left|\frac{u(w)}{(1-\overline{z}\varphi(w))^{2+\alpha}}-
\frac{v(w)}{(1-\overline{z}\psi(w))^{2+\alpha}}\right|^2(1-|z|^2)^{2+\alpha}{\rm d}A_{\alpha}(w)\\
&\gtrsim \frac{\int_{\varphi^{-1}(D(z,r))}\left|u(w)-v(w)\frac{(1-\overline{z}\varphi(w))^{2+\alpha}}{(1-\overline{z}\psi(w))^{2+\alpha}}\right|^2{\rm d}A_{\alpha}(w)}{(1-|z|^2)^{2+\alpha}}.
\end{split}
\end{equation}
Similarly,
\begin{equation}\label{equa4.5}
\begin{split}
&\|(C_{u,\varphi}-C_{v,\psi})k_z^{[1]}\|_{A_{\alpha}^2}^2\\
&\gtrsim\frac{\int_{\varphi^{-1}(D(z,r))}\left|u(w)\varphi(w)-v(w)\psi(w)\frac{(1-\overline{z}\varphi(w))^{3+\alpha}}{(1-\overline{z}\psi(w))^{3+\alpha}}\right|^2{\rm d}A_{\alpha}(w)}{(1-|z|^2)^{2+\alpha}}.
\end{split}
\end{equation}
Since $\left|\frac{1-\overline{z}\varphi(w)}{1-\overline{z}\psi(w)}\right|\lesssim 1$ when $w\in \varphi^{-1}(D(z,r))$, multiplying the integrand in \eqref{equa4.4} by $|\psi(w)|^2\left|\frac{1-\overline{z}\varphi(w)}{1-\overline{z}\psi(w)}\right|^2$, and adding it to \eqref{equa4.5}, by the triangle inequality and \eqref{equa2.4}, we obtain
\begin{equation}\label{equa4.6}
\begin{split}
&\|(C_{u,\varphi}-C_{v,\psi})k_z^{[0]}\|_{A_{\alpha}^2}^2+\|(C_{u,\varphi}-C_{v,\psi})k_z^{[1]}\|_{A_{\alpha}^2}^2\\
&\quad \gtrsim\frac{\int_{\varphi^{-1}(D(z,r))}|u(w)|^2\rho(w)^2{\rm d}A_{\alpha}(w)}{(1-|z|^2)^{2+\alpha}}=M_r(\omega_{\varphi,u}^2)(z).
\end{split}
\end{equation}
Thus $M_r(\omega_{\varphi,u}^2)\in L^{\frac{p}{2}}(\mathbb{D},{\rm d}\lambda)$. Similarly, $M_r(\omega_{\psi,v}^2)\in L^{\frac{p}{2}}(\mathbb{D},{\rm d}\lambda)$.

On the other hand, when $\varphi(w)\in D(z,r)$, by \eqref{equa2.4}, we have
\begin{equation*}
\begin{split}
&\left|1-\frac{(1-\overline{z}\varphi(w))^{\alpha+2}}{(1-\overline{z}\psi(w))^{\alpha+2}}\right|\\
&\quad\lesssim \sup_{0\leq t\leq 1}\frac{|1-\overline{z}(t\varphi(w)+(1-t)\psi(w))|^{\alpha+1}}{|1-\overline{z}\psi(w)|^{\alpha+1}}\frac{|\varphi(w)-\psi(w)|}{|1-\overline{z}\psi(w)|}\lesssim \rho(z,w).
\end{split}
\end{equation*}
It follows that when $\varphi(w)\in D(z,r)$, 
\begin{align*}
&|u(w)-v(w)|(1-\rho(w))^{2+\alpha}\\
&\quad\lesssim |u(w)-v(w)|\left|\frac{1-\overline{z}\varphi(w)}{1-\overline{z}\psi(w)}\right|^{\alpha+2}\\
&\quad \leq \left|u(w)-v(w)\frac{(1-\overline{z}\varphi(w))^{\alpha+2}}{(1-\overline{z}\psi(w))^{\alpha+2}}\right|+|u(w)|\left|1-\frac{(1-\overline{z}\varphi(w))^{\alpha+2}}{(1-\overline{z}\psi(w))^{\alpha+2}}\right|\\
&\quad \lesssim \left|u(w)-v(w)\frac{(1-\overline{z}\varphi(w))^{\alpha+2}}{(1-\overline{z}\psi(w))^{\alpha+2}}\right|+|u(w)|\rho(w).
\end{align*}
So when $\beta\geq 2(\alpha+2)$ we combine \eqref{equa4.4} and \eqref{equa4.6} to obtain
\begin{equation*}
\begin{split}
&\|(C_{u,\varphi}-C_{v,\psi})k_z^{[0]}\|_{A_{\alpha}^2}^2+\|(C_{u,\varphi}-C_{v,\psi})k_z^{[1]}\|_{A_{\alpha}^2}^2\\
&\quad \gtrsim\frac{\int_{\varphi^{-1}(D(z,r))}|u(w)-v(w)|^2(1-\rho(w))^{2(2+\alpha)}{\rm d}A_{\alpha}(w)}{(1-|z|^2)^{2+\alpha}}\geq M_r(\sigma_{\varphi,\beta}^2)(z).
\end{split}
\end{equation*}
Consequently, $M_r(\sigma_{\varphi,\beta}^2)\in L^{\frac{p}{2}}(\mathbb{D},{\rm d}\lambda)$. Similarly, $M_r(\sigma_{\psi,\beta}^2)\in L^{\frac{p}{2}}(\mathbb{D},{\rm d}\lambda)$. The proof is complete.
\end{proof}

\begin{remark}\label{remark4.1}
According to Theorem \ref{theorem1.3}, $C_{u,\varphi}-C_{v,\psi}$ is Hilbert-Schmidt on $A_{\alpha}^2(\mathbb{D})$ if and only if
\begin{itemize}
\item[(a)] $M_r(\omega_{2})$ and $M_r(\sigma_{\beta,2})$ belong to $L^1(\mathbb{D},{\rm d}\lambda)$ for some (or any) $r>0$. Here, $\beta=2(2+\alpha)$.
\end{itemize}
This is equivalent to \cite[Theorem 2]{1}, i.e. $C_{u,\varphi}-C_{v,\psi}$ is Hilbert-Schmidt on $A_{\alpha}^2(\mathbb{D})$ if and only if 
\begin{itemize}
\item[(b)] the functions $\frac{|\rho u|^2}{(1-|\varphi|^2)^{2+\alpha}}$, $\frac{|\rho v|^2}{(1-|\psi|^2)^{2+\alpha}}$, $\frac{(1-\rho)^{2+\alpha}|u-v|^2}{(1-|\varphi|^2)^{2+\alpha}}$ and $\frac{(1-\rho)^{2+\alpha}|u-v|^2}{(1-|\psi|^2)^{2+\alpha}}$ belong to $L^1(\mathbb{D},{\rm d}A_{\alpha})$.
\end{itemize}

In fact, by \eqref{equa2.3} and Fubini's Theorem, 
\begin{equation}\label{equa4.7}
\begin{split}
&\int_{\mathbb{D}}M_r(\omega_2)(z){\rm d}\lambda(z)\\
&\quad =\int_{\mathbb{D}}\frac{\int_{\varphi^{-1}(D(z,r))}|u(w)|^2\rho(w)^2{\rm d}A_{\alpha}(w)}{(1-|z|^2)^{2+\alpha}}{\rm d}\lambda(z)\\
&\qquad \quad +\int_{\mathbb{D}}\frac{\int_{\psi^{-1}(D(z,r))}|v(w)|^2\rho(w)^2{\rm d}A_{\alpha}(w)}{(1-|z|^2)^{2+\alpha}}{\rm d}\lambda(z)\\
&\quad =\int_{\mathbb{D}}|u(w)|^2\rho(w)^2\left(\int_{D(\varphi(w),r)}\frac{1}{(1-|z|^2)^{2+\alpha}}{\rm d}\lambda(z)\right){\rm d}A_{\alpha}(w)\\
&\qquad \quad +\int_{\mathbb{D}}|v(w)|^2\rho(w)^2\left(\int_{D(\psi(w),r)}\frac{1}{(1-|z|^2)^{2+\alpha}}{\rm d}\lambda(z)\right){\rm d}A_{\alpha}(w)\\
&\quad \simeq \int_{\mathbb{D}}\frac{|u(w)|^2\rho(w)^2}{(1-|\varphi(w)|^2)^{2+\alpha}}+\frac{|v(w)|^2\rho(w)^2}{(1-|\psi(w)|^2)^{2+\alpha}}{\rm d}A_{\alpha}(w)
\end{split}
\end{equation}
Meanwhile, for $\beta=2(\alpha+2)$,
\begin{equation}\label{equa4.8}
\begin{split}
&\int_{\mathbb{D}}M_r(\sigma_{\beta,2}){\rm d}\lambda\\
&\quad \simeq\int_{\mathbb{D}}|u-v|^2(1-\rho)^{2(2+\alpha)}\left(\frac{1}{(1-|\varphi|^2)^{2+\alpha}}+\frac{1}{(1-|\psi|^2)^{2+\alpha}}\right){\rm d}A_{\alpha}\\
&\quad \leq\int_{\mathbb{D}}|u-v|^2(1-\rho)^{2+\alpha}\left(\frac{1}{(1-|\varphi|^2)^{2+\alpha}}+\frac{1}{(1-|\psi|^2)^{2+\alpha}}\right){\rm d}A_{\alpha}.
\end{split}
\end{equation}
On the other hand, for fixed $0<\delta<1$, let $E_{\delta}=\{z:\rho(z)<\delta\}$. We write
\begin{equation*}
\begin{split}
&\int_{\mathbb{D}}|u-v|^2(1-\rho)^{2+\alpha}\left(\frac{1}{(1-|\varphi|^2)^{2+\alpha}}+\frac{1}{(1-|\psi|^2)^{2+\alpha}}\right){\rm d}A_{\alpha}\\
&\quad =\int_{E_{\delta}}+\int_{\mathbb{D}\backslash E_{\delta}}|u-v|^2(1-\rho)^{2+\alpha}\left(\frac{1}{(1-|\varphi|^2)^{2+\alpha}}+\frac{1}{(1-|\psi|^2)^{2+\alpha}}\right){\rm d}A_{\alpha}.
\end{split}
\end{equation*}
Clearly, 
\begin{equation}\label{equa4.9}
\begin{split}
\int_{E_{\delta}}|u-v|^2(1-\rho)^{2+\alpha}\left(\frac{1}{(1-|\varphi|^2)^{2+\alpha}}+\frac{1}{(1-|\psi|^2)^{2+\alpha}}\right){\rm d}A_{\alpha} \lesssim \int_{\mathbb{D}}M_{r}(\sigma_{\beta,2}){\rm d}\lambda.
\end{split}
\end{equation}
And by \eqref{equa2.1}, we have
\begin{equation}\label{equa4.10}
\begin{split}
&\int_{\mathbb{D}\backslash E_{\delta}}|u-v|^2(1-\rho)^{2+\alpha}\left(\frac{1}{(1-|\varphi|^2)^{2+\alpha}}+\frac{1}{(1-|\psi|^2)^{2+\alpha}}\right){\rm d}A_{\alpha}\\
&\quad =\int_{\mathbb{D}\backslash E_{\delta}}|u-v|^2\frac{(1-|\varphi|^2)^{2+\alpha}+(1-|\psi|^2)^{2+\alpha}}{|1-\overline{\varphi}\psi|^{2(2+\alpha)}}{\rm d}A_{\alpha}\\
&\quad \lesssim \int_{\mathbb{D}\backslash E_{\delta}}\frac{(|u|^2+|v|^2)\rho^2}{|1-\overline{\varphi}\psi|^{2+\alpha}}{\rm d}A_{\alpha}\lesssim \int_{\mathbb{D}}M_r(\omega_2){\rm d}\lambda.
\end{split}
\end{equation}
Combining \eqref{equa4.7}-\eqref{equa4.10}, the equivalence of (a) and (b) holds.
\end{remark}

From the proof of Theorem \ref{theorem1.2}, we conclude the following characterization for Schatten class difference $C_{u,\varphi}-C_{v,\psi}$ via Reproducing Kernel Thesis.

\begin{corollary}\label{corollary4.2}
Let $\alpha>-1$ and $p\geq 2$. Suppose $u,v\in H(\mathbb{D})$ and $\varphi,\psi\in S(\mathbb{D})$. Then $C_{u,\varphi}-C_{v,\psi}\in S_p(A_{\alpha}^2)$ if and only if 
\[\int_{\mathbb{D}}\|(C_{u,\varphi}-C_{v,\psi})k_z^{[n]}\|_{A_{\alpha}^2}^p{\rm d}\lambda(z)<\infty\]
for any $n\in\mathbb{N}$. 
\end{corollary}

\begin{remark}\label{remark3.3}
Let $v(z)=0$, when $p\geq 2$, by Theorem \ref{theorem1.3} and Corollary \ref{corollary4.2}, it is easy to see that $C_{u,\varphi}\in S_p(A_{\alpha}^2)$ if and only if
\begin{equation}\label{equa4.11}
\int_{\mathbb{D}}\left(\frac{\int_{\varphi^{-1}(D(z,r))}|u(w)|^2{\rm d}A_{\alpha}(w)}{(1-|z|^2)^{2+\alpha}}\right)^{\frac{p}{2}}{\rm d}\lambda(z)<\infty,
\end{equation}
if and only if
\begin{equation*}
\int_{\mathbb{D}}\|C_{u,\varphi}k_z^{[n]}\|_{A_{\alpha}^2}^p{\rm d}\lambda(z)<\infty,\quad \forall n\in\mathbb{N}.
\end{equation*}
In fact, an easy application of \eqref{equa2.8} shows that \eqref{equa4.11} holds for all $0<p<\infty$, since $C_{u,\varphi}^*C_{u,\varphi}=T_{\mu_{u,\varphi}}$, where $\mu_{u,\varphi}=(|u|^2dA_{\alpha})\circ\varphi^{-1}$.
\end{remark}

\begin{proposition}\label{proposition4.4}
Let $\alpha>-1$, $p\geq 2$ and $a,b\in\mathbb{C}\backslash\{0\}$. Suppose $C_{\varphi}$ and $C_{\psi}$ are not in $S_{p}(A_{\alpha}^2)$. Then $aC_{\varphi}+bC_{\psi}\in S_p(A_{\alpha}^2)$ if and only if $a+b=0$ and $C_{\varphi}-C_{\psi}\in S_p(A_{\alpha}^2)$.
\end{proposition}

\begin{proof}
The sufficiency is trivial. We only need to prove the necessity.

Assume $aC_{\varphi}+bC_{\psi}\in S_p(A_{\alpha}^2)$. Fix $r>0$. By Theorem \ref{theorem1.3}, we have
\begin{equation}\label{equa4.12}
\int_{\mathbb{D}}\left(\frac{\int_{\varphi^{-1}(D(z,r))}\rho(w)^2{\rm d}A_{\alpha}(w)}{(1-|z|^2)^{2+\alpha}}\right)^{\frac{p}{2}}{\rm d}\lambda(z)<\infty,
\end{equation} 
and
\begin{equation}\label{equa4.13}
\int_{\mathbb{D}}\left(\frac{\int_{\varphi^{-1}(D(z,r))}(1-\rho(w))^{2(2+\alpha)}|a+b|^2{\rm d}A_{\alpha}(w)}{(1-|z|^2)^{2+\alpha}}\right)^{\frac{p}{2}}{\rm d}\lambda(z)<\infty.
\end{equation}
If $a+b\neq 0$, then adding \eqref{equa4.12} and \eqref{equa4.13}, we easily obtain
\[\int_{\mathbb{D}}\left(\frac{A_{\alpha}(\varphi^{-1}(D(z,r)))}{(1-|z|^2)^{\alpha+2}}\right)^{\frac{p}{2}}{\rm d}\lambda(z)<\infty.\]
By \eqref{equa4.11}, this shows that $C_{\varphi}\in S_p(A_{\alpha}^2)$, which is a contradiction. It follows that $a+b=0$ and $C_{\varphi}-C_{\psi}\in S_p(A_{\alpha}^2)$.
\end{proof}



\section{Proof of Theorem 1.4}

In this section, we aim to prove Theoremm \ref{theorem1.4}. To this end, we need the following lemma.

\begin{lemma}\label{lemma5.1}
Let $u\in H(\mathbb{D})$ and $\varphi\in S(\mathbb{D})$. Then $C_{u,\varphi}$ is bounded (compact, resp.) on $H^2(\mathbb{D})$ if and only if both $C_{u',\varphi}: H^2(\mathbb{D})\to A_1^2(\mathbb{D})$ and $C_{u\varphi',\varphi}: A_1^2(\mathbb{D})\to A_{1}^2(\mathbb{D})$ are bounded (compact, resp.).
\end{lemma}

\begin{proof}
By \eqref{equa1.1}, we know that $f\in H^2(\mathbb{D})$ if and only if $f'\in A_1^2(\mathbb{D})$. So the suffciency follows from the fact that 
\begin{align*}
\|(C_{u,\varphi}f)'\|_{A_1^2}^2&=\int_{\mathbb{D}}|u'(z)f(\varphi(z))+u(z)\varphi'(z)f'(\varphi(z))|^2{\rm d}A_1(z)\\
&\lesssim \|C_{u',\varphi}f\|_{A_1^2}^2+\|C_{u\varphi',\varphi}f'\|_{A_{1}^2}^2.
\end{align*}

Now we prove the necessity. Assume $C_{u,\varphi}$ is bounded on $H^2(\mathbb{D})$, we use the test function $g_{a,N,2}^{[i]}$ defined in Section 2.2 to obtain
$$\sup_{a\in\mathbb{D}}\|C_{u,\varphi}g_{a,N,2}^{[i]}\|_{H^2}<\infty,\quad i\in\mathbb{N}.$$
By \eqref{equa1.1}, taking $i=0$, we have
\begin{align*}
&\|C_{u,\varphi}g_{a,N,2}^{[0]}\|_{H^2}^2\\
&\quad \gtrsim \int_{\mathbb{D}}|(C_{u,\varphi}g_{a,N,2}^{[0]})'(z)|^2{\rm d}A_1(z)\\
&\quad =\int_{\mathbb{D}}\left|\frac{u'(z)}{(1-\overline{a}\varphi(z))^N}+\frac{u(z)\varphi'(z)N\overline{a}}{(1-\overline{a}\varphi(z))^{N+1}}\right|^2(1-|a|^2)^{2N-1}{\rm d}A_1(z).
\end{align*}
So we get
\begin{align}\label{equa5.1}
\sup_{a\in\mathbb{D}}\int_{\mathbb{D}}\left|\frac{u'(z)}{(1-\overline{a}\varphi(z))^N}+\frac{u(z)\varphi'(z)N\overline{a}}{(1-\overline{a}\varphi(z))^{N+1}}\right|^2(1-|a|^2)^{2N-1}{\rm d}A_1(z)<\infty.
\end{align}
Similarly, taking $i=1$, we get
\begin{align}\label{equa5.2}
\sup_{a\in\mathbb{D}}\int_{\mathbb{D}}\left|\frac{u'(z)\varphi(z)}{(1-\overline{a}\varphi(z))^{N+1}}+\frac{u(z)\varphi'(z)(1+N\overline{a}\varphi(z))}{(1-\overline{a}\varphi(z))^{N+2}}\right|^2(1-|a|^2)^{2N+1}{\rm d}A_1(z)<\infty.
\end{align}
Note that $\left|\frac{1-|a|^2}{1-\overline{a}\varphi(z)}\right|\lesssim 1$. Multiplying the integrand in \eqref{equa5.1} by $\frac{(1-|a|^2)^2}{|1-\overline{a}\varphi(z)|^2}|\varphi(z)|^2$, adding it to \eqref{equa5.2} and by the triangle inequality, we obtain
\begin{align}\label{equa5.3}
\sup_{a\in\mathbb{D}}\int_{\mathbb{D}}|u(z)\varphi'(z)|^2\frac{(1-|a|^2)^{2N+1}}{|1-\overline{a}\varphi(z)|^{2N+4}}{\rm d}A_1(z)<\infty.
\end{align}
Then it follows from \eqref{equa2.3} that
\begin{align}\label{equa5.4}
\sup_{a\in\mathbb{D}}\frac{\int_{\varphi^{-1}(D(a,r))}|u(z)\varphi'(z)|^2{\rm d}A_1(z)}{(1-|a|^2)^3}<\infty,
\end{align}
which means that $(|u\varphi'|{\rm d}A_1)\circ\varphi^{-1}$ is a Carleson measure for $A_1^2(\mathbb{D})$ by Lemma \ref{lemma2.7}. So $C_{u\varphi',\varphi}$ is bounded on $A_1^2(\mathbb{D})$.

Since \eqref{equa5.1} also holds replacing $N$ by $N+1$, combining \eqref{equa5.1} and \eqref{equa5.3}, we get
\begin{align}\label{equa5.5}
\sup_{a\in\mathbb{D}}\int_{\mathbb{D}}|u'(z)|^2\frac{(1-|a|^2)^{2N+1}}{|1-\overline{a}\varphi(z)|^{2N+2}}{\rm d}A_1(z)<\infty.
\end{align}
This shows that $(|u'|{\rm d}A_1)\circ\varphi^{-1}$ is a Hardy Carleson measure by Equation (1.1) in \cite{25}. It follows that $C_{u',\varphi}: H^2(\mathbb{D})\to A_1^2(\mathbb{D})$ is bounded.

The proof for the compactness part is just a modification. We omit the routine details.
\end{proof}

We are now ready to prove Theorem \ref{theorem1.4}.

\begin{proof}[{\bf Proof of Theorem 1.4}]
The sufficiency is trivial, since for $f\in H^2(\mathbb{D})$, we have $f'\in A_1^2(\mathbb{D})$ and 
\begin{align*}
\left\|\left[(C_{u,\varphi}-C_{v,\psi})f\right]'\right\|_{A_1^2}\lesssim \|(C_{u',\varphi}-C_{v',\psi})f\|_{A_1^2}+\|(C_{u\varphi',\varphi}-C_{v\psi',\psi})f'\|_{A_1^2}.
\end{align*}

Now we prove the necessity. Assume both $C_{u,\varphi}$ and $C_{v,\psi}$ are bounded on $H^2(\mathbb{D})$ and $C_{u,\varphi}-C_{v,\psi}$ is compact on $H^2(\mathbb{D})$. By the triangle inequality, for $f\in H^2(\mathbb{D})$, we have
\begin{align*}
\|(C_{u',\varphi}-C_{v',\psi})f\|_{A_1^2}\lesssim \|(C_{u,\varphi}-C_{v,\psi})f\|_{H^2}+\|(C_{u\varphi',\varphi}-C_{v\psi',\psi})f'\|_{A_1^2}.
\end{align*}
Thus we only need to prove the compactness of $C_{u\varphi',\varphi}-C_{v\psi',\psi}$ on $A_1^2(\mathbb{D})$.

Recall the test function 
$$g_{a,N,2}^{[i]}(z)=\frac{(1-|a|^2)^{N+i-\frac{1}{2}}z^i}{(1-\overline{a}z)^{N+i}},\quad z\in\mathbb{D}.$$
By Lemma \ref{lemma2.4}, the compactness of $C_{u,\varphi}-C_{v,\psi}$ on $H^2(\mathbb{D})$ implies that
\begin{align*}
\lim_{|a|\to 1}\|(C_{u,\varphi}-C_{v,\psi})g_{a,N,2}^{[i]}\|_{H^2}=0,\quad i\in\mathbb{N}.
\end{align*}
By \eqref{equa1.1}, taking $i=0$, we have
\begin{align*}
&\|(C_{u,\varphi}-C_{v,\psi})g_{a,N,2}^{[0]}\|_{H^2}^2\\
&\quad \gtrsim \int_{\mathbb{D}}\left|\left[(C_{u,\varphi}-C_{v,\psi})g_{a,N,2}^{[0]}\right]'(z)\right|^2{\rm d}A_1(z)\\
&\quad =\int_{\mathbb{D}}\left|\frac{u'(z)}{(1-\overline{a}\varphi(z))^{N}}+\frac{u(z)\varphi'(z)N\overline{a}}{(1-\overline{a}\varphi(z))^{N+1}}\right.\\
&\quad\qquad\qquad \left.-\frac{v'(z)}{(1-\overline{a}\psi(z))^{N}}-\frac{v(z)\psi'(z)N\overline{a}}{(1-\overline{a}\psi(z))^{N+1}}\right|^2(1-|a|^2)^{2N-1}{\rm d}A_{1}(z)
\end{align*}
Fix $r>0$. Then by \eqref{equa2.3}, we have
\begin{align*}
&\|(C_{u,\varphi}-C_{v,\psi})g_{a,N,2}^{[0]}\|_{H^2}^2\\
&\quad \gtrsim\frac{\int_{\varphi^{-1}(D(a,r))}\left|\left(u'(z)+\frac{u(z)\varphi'(z)N\overline{a}}{1-\overline{a}\varphi(z)}\right)-\left(v'(z)+\frac{v(z)\psi'(z)N\overline{a}}{1-\overline{a}\psi(z)}\right)\left(\frac{1-\overline{a}\varphi(z)}{1-\overline{a}\psi(z)}\right)^N\right|^2{\rm d}A_1(z)}{(1-|a|^2)}.
\end{align*}
It follows that
\begin{align}\label{equa5.6}
&\lim_{|a|\to 1}\int\limits_{\varphi^{-1}(D(a,r))}\left|\left(v'(z)+\frac{v(z)\psi'(z)N\overline{a}}{1-\overline{a}\psi(z)}\right)\left(\frac{1-\overline{a}\varphi(z)}{1-\overline{a}\psi(z)}\right)^N\right.\nonumber \\
&\qquad\qquad\qquad\qquad\left.-\left(u'(z)+\frac{u(z)\varphi'(z)N\overline{a}}{1-\overline{a}\varphi(z)}\right)\right|^2\frac{{\rm d}A_1(z)}{(1-|a|^2)}=0
\end{align}
Similarly, taking $i=1$ and $i=2$, respectively, we obtain
\begin{align}\label{equa5.7}
&\lim_{|a|\to 1}\int\limits_{\varphi^{-1}(D(a,r))}\left|\left(v'(z)\psi(z)+\frac{v(z)\psi'(z)(1+N\overline{a}\psi(z))}{1-\overline{a}\psi(z)}\right)\left(\frac{1-\overline{a}\varphi(z)}{1-\overline{a}\psi(z)}\right)^{N+1}\right.\nonumber\\
&\qquad\qquad\qquad\left.-\left(u'(z)\varphi(z)+\frac{u(z)\varphi'(z)(1+N\overline{a}\varphi(z))}{1-\overline{a}\varphi(z)}\right)\right|^2\frac{{\rm d}A_1(z)}{1-|a|^2}=0,
\end{align}
and 
\begin{align}\label{equa5.8}
&\lim_{|a|\to 1}\int\limits_{\varphi^{-1}(D(a,r))}\left|\left(v'(z)\psi(z)^2+\frac{v(z)\psi(z)(2\psi(z)+N\overline{a}\psi(z)^2)}{1-\overline{a}\psi(z)}\right)\left(\frac{1-\overline{a}\varphi(z)}{1-\overline{a}\psi(z)}\right)^{N+2}\right.\nonumber\\
&\qquad\qquad\left.-\left(u'(z)\varphi(z)^2+\frac{u(z)\varphi'(z)(2\varphi(z)+N\overline{a}\varphi(z)^2)}{1-\overline{a}\varphi(z)}\right)\right|^2\frac{{\rm d}A_1(z)}{1-|a|^2}=0.
\end{align}

By \eqref{equa2.3}, $\left|\frac{1-\overline{a}\varphi(z)}{1-\overline{a}\psi(z)}\right|\lesssim 1$ when $\varphi(z)\in D(a,r)$. Multiplying the integrand in \eqref{equa5.6} by $\left|\frac{1-\overline{a}\varphi(z)}{1-\overline{a}\psi(z)}\right|^2|\psi(z)|^2$, adding it to \eqref{equa5.7} and by the triangle inequality, we obtain
\begin{align}\label{equa5.9}
&\lim_{|a|\to 1}\int\limits_{\varphi^{-1}(D(a,r))}\left|\frac{v(z)\psi'(z)}{1-\overline{a}\psi(z)}\left(\frac{1-\overline{a}\varphi(z)}{1-\overline{a}\psi(z)}\right)^{N+1}-\frac{u(z)\varphi'(z)}{1-\overline{a}\varphi(z)}\right.\nonumber\\
&\qquad\qquad\left.-\left(u'(z)+\frac{u(z)\varphi'(z)N\overline{a}}{1-\overline{a}\varphi(z)}\right)\frac{\varphi(z)-\psi(z)}{1-\overline{a}\psi(z)}\right|^2\frac{{\rm d}A_1(z)}{1-|a|^2}=0.
\end{align}
Similarly, multiplying the integrand in \eqref{equa5.7} by $\left|\frac{1-\overline{a}\varphi(z)}{1-\overline{a}\psi(z)}\right|^2|\psi(z)|^2$, adding it to \eqref{equa5.8} and by the triangle inequality, we obtain
\begin{align}\label{equa5.10}
&\lim_{|a|\to 1}\int\limits_{\varphi^{-1}(D(a,r))}\left|\frac{v(z)\psi'(z)\psi(z)}{1-\overline{a}\psi(z)}\left(\frac{1-\overline{a}\varphi(z)}{1-\overline{a}\psi(z)}\right)^{N+2}-\frac{u(z)\varphi'(z)\varphi(z)}{1-\overline{a}\varphi(z)}\right.\nonumber\\
&\qquad\qquad\left.-\left(u'(z)\varphi(z)+\frac{u(z)\varphi'(z)(1+N\overline{a}\varphi(z))}{1-\overline{a}\varphi(z)}\right)\frac{\varphi(z)-\psi(z)}{1-\overline{a}\psi(z)}\right|^2\frac{{\rm d}A_1(z)}{1-|a|^2}=0.
\end{align}
Again, multiplying the integrand in \eqref{equa5.9} by $\left|\frac{1-\overline{a}\varphi(z)}{1-\overline{a}\psi(z)}\right|^2|\psi(z)|^2$, adding it to \eqref{equa5.10} and by the triangle inequality, we obtain
\begin{align}\label{equa5.11}
&\lim_{|a|\to 1}\int\limits_{\varphi^{-1}(D(a,r))}\left|2\frac{u(z)\varphi'(z)}{1-\overline{a}\varphi(z)}\frac{\varphi(z)-\psi(z)}{1-\overline{a}\psi(z)}\right.\nonumber\\
&\qquad\qquad\left.+\left(u'(z)+\frac{u(z)\varphi'(z)N\overline{a}}{1-\overline{a}\varphi(z)}\right)\left(\frac{\varphi(z)-\psi(z)}{1-\overline{a}\psi(z)}\right)^2\right|^2\frac{{\rm d}A_1(z)}{1-|a|^2}=0.
\end{align}
Since \eqref{equa5.11} also holds replacing $N$ by $N+1$, we get 
\begin{align*}
\lim_{|a|\to 1}\int_{\varphi^{-1}(D(a,r))}\left|\frac{u(z)\varphi'(z)}{1-\overline{a}\varphi(z)}\right|^2\left|\frac{\varphi(z)-\psi(z)}{1-\overline{a}\psi(z)}\right|^4\frac{{\rm d}A_1(z)}{1-|a|^2}=0.
\end{align*}
Then it follows from \eqref{equa2.3} and \eqref{equa2.4} that
\begin{align}\label{equa5.12}
\lim_{|a|\to 1}\frac{\int_{\varphi^{-1}(D(a,r))}|u(z)\varphi'(z)|^2\rho(z)^4{\rm d}A_1(z)}{(1-|a|^2)^3}=0.
\end{align}
Since $C_{u,\varphi}$ is bounded on $H^2(\mathbb{D})$, combining \eqref{equa5.12} and \eqref{equa5.3}, we obtain
\begin{align}\label{equa5.13}
\lim_{|a|\to 1}\frac{\int_{\varphi^{-1}(D(a,r))}|u(z)\varphi'(z)|^2\rho(z)^2{\rm d}A_1(z)}{(1-|a|^2)^3}=0.
\end{align}
This shows that $(|u\varphi'|^2\rho^2{\rm d}A_1)\circ\varphi^{-1}$ is a vanishing Carleson measure for $A_1^2(\mathbb{D})$. Similarly, $(|v\psi'|^2\rho^2{\rm d}A_1)\circ\psi^{-1}$ is also a vanishing Carleson measure for $A_1^2(\mathbb{D})$.

On the other hand, using \eqref{equa2.3} and \eqref{equa2.4}, combining \eqref{equa5.13} and \eqref{equa5.11}, we obtain
\begin{align}\label{equa5.14}
\lim_{|a|\to 1}\frac{\int_{\varphi^{-1}(D(a,r))}|u'(z)|^2\rho(z)^4{\rm d}A_1(z)}{1-|a|^2}=0.
\end{align}
Since $C_{u,\varphi}$ is bounded on $H^2(\mathbb{D})$, by \eqref{equa5.5} and \eqref{equa2.3}, we have
\begin{align*}
\sup_{a\in\mathbb{D}}\frac{\int_{\varphi^{-1}(D(a,r))}|u'(z)|^2{\rm d}A_1(z)}{1-|a|^2}<\infty.
\end{align*}
This, together with \eqref{equa5.14}, implies that
\begin{align}\label{equa5.15}
\lim_{|a|\to 1}\frac{\int_{\varphi^{-1}(D(a,r))}|u'(z)|^2\rho(z)^2{\rm d}A_1(z)}{1-|a|^2}=0.
\end{align}
Now using \eqref{equa2.3} and \eqref{equa2.4} again, we combine \eqref{equa5.15}, \eqref{equa5.13} and \eqref{equa5.9} to obtain
\begin{align*}
\lim_{|a|\to 1}\frac{\int_{\varphi^{-1}(D(a,r))}\left|v(z)\psi'(z)\left(\frac{1-\overline{a}\varphi(z)}{1-\overline{a}\psi(z)}\right)^{N+2}-u(z)\varphi'(z)\right|^2{\rm d}A_1(z)}{(1-|a|^2)^3}=0.
\end{align*}
Then modifying the proof of $(i)\Rightarrow(ii)$ in Theorem \ref{theorem1.1} (see \eqref{equa3.2}), we obtain $(|u\varphi'-v\psi'|^2(1-\rho)^{2(N+2)}{\rm d}A_1)\circ\varphi^{-1}$ is a vanishing Carleson measure for $A_1^2(\mathbb{D})$. Similarly, $(|u\varphi'-v\psi'|^2(1-\rho)^{2(N+2)}{\rm d}A_1)\circ\psi^{-1}$ is also a vanishing Carleson measure for $A_1^2(\mathbb{D})$. Finally, according to Theorem \ref{theorem1.1}, $C_{u\varphi',\varphi}-C_{v\psi',\psi}$ is compact on $A_1^2(\mathbb{D})$. The proof is complete.
\end{proof}


\begin{theorem}\label{theorem5.2}
Let $a,b\in \mathbb{C}\backslash\{0\}$ and $\varphi,\psi\in S(\mathbb{D})$ be of bounded valence. Suppose $C_{\varphi}$ and $C_{\psi}$ are not compact on $H^2(\mathbb{D})$. Then $aC_{\varphi}+bC_{\psi}$ is compact on $H^2(\mathbb{D})$ if and only if $a+b=0$ and $C_{\varphi}-C_{\psi}$ is compact on $H^2(\mathbb{D})$.
\end{theorem}

\begin{proof}
The sufficiency is trivial. We only need to prove the necessity.

Assume $aC_{\varphi}+bC_{\psi}$ is compact on $H^2(\mathbb{D})$. Then $aC_{\varphi',\varphi}-bC_{\psi',\psi}$ is compact on $A_1^2(\mathbb{D})$. Let $r>0$ and $\beta>3$. By Theorem \ref{theorem1.1}, we have
\begin{align}\label{equa5.16}
\lim_{|w|\to 1}\frac{\int_{\varphi^{-1}(D(w,r))}\rho(z)^2|\varphi'(z)|^2{\rm d}A_1(z)}{(1-|w|^2)^3}=0,
\end{align}
and 
\begin{align}\label{equa5.17}
\lim_{|w|\to 1}\frac{\int_{\varphi^{-1}(D(w,r))}(1-\rho(z))^{\beta}|a\varphi'(z)+b\psi'(z)|^2{\rm d}A_1(z)}{(1-|w|^2)^3}=0.
\end{align}
For $\delta\in (0,1)$, recall that $E_{\delta}=\{z: \rho(z)<\delta\}$. By the proof of \cite[Theorem 1.3]{4}, if $\varphi$ and $\psi$ are of bounded valence, then
\begin{align}\label{equa5.18}
\lim_{|w|\to 1}\frac{\int_{\varphi^{-1}(D(w,r))\cap E_{\delta}}|\varphi'(z)-\psi'(z)|^2{\rm d}A_1(z)}{(1-|w|^2)^3}=0.
\end{align}
Note that $1-\rho(z)\geq 1-\delta$ for $z\in E_\delta$. If $a+b\neq 0$, then by \eqref{equa5.17} and \eqref{equa5.18}, we get
\begin{align*}
\lim_{|w|\to 1}\frac{\int_{\varphi^{-1}(D(w,r))\cap E_{\delta}}|\varphi'(z)|^2{\rm d}A_1(z)}{(1-|w|^2)^3}=0.
\end{align*}
This, together with \eqref{equa5.16}, implies that
\begin{align*}
\lim_{|w|\to 1}\frac{\int_{\varphi^{-1}(D(w,r))}|\varphi'(z)|^2{\rm d}A_1(z)}{(1-|w|^2)^3}=0.
\end{align*}
Then by Lemma \ref{lemma2.7}, $(|\varphi'|^2{\rm d}A_1)\circ\varphi^{-1}$ is a compact Bergman Carleson measure. It follows that $C_{\varphi',\varphi}$ is compact on $A_1^2(\mathbb{D})$. Then $C_{\varphi}$ is compact on $H^2(\mathbb{D})$, which is a contradiction. So we have $a+b=0$ and $C_{\varphi}-C_{\psi}$ is compact on $H^2(\mathbb{D})$.
\end{proof}

% ------------------------------------------------------------------------


\noindent{\bf Acknowledgment}
Tong was supported in part by the National Natural Science Foundation of China (Grant No. 12171136, 12411530045). Yang was supported in part by the National Natural Science Foundation of China (Grant No. 12501103) and the Natural Science Foundation of Hebei Province (Grant No. A2023202031, A2023202037).\\

\noindent{\bf Declaration}
The authors declare that they have no conflicts of interest. No data was used for the research described in this article. 


\begin{thebibliography}{99}
	
	\bibitem{1} \textsc{S. Acharyya, Z. Wu}, Compact and Hilbert-Schmidt differences of weighted composition operators, \textit{Integral. Equ. Oper Theory}, \textbf{88} (2017), 465-482.

\bibitem{2} \textsc{E. Berkon}, Composition operators isolated in the uniform operator topology, \textit{Proc. Amer. Math. Soc.}, \textbf{81} (1981), 230-232.

\bibitem{3} \textsc{B. Choe, K. Choi, H. Koo, I. Park}, Difference of weighted composition operators over the ball, \textit{Complex Anal. Oper. Theory}, \textbf{18} (2024), Paper No. 33, 34pp.

\bibitem{4} \textsc{B. Choe, K. Choi, H. Koo, J. Yang}, Difference of weighted composition operators, \textit{J. Funct. Anal.}, \textbf{278} (2020), 108401, 38pp.

\bibitem{5} \textsc{B. Choe, K. Choi, H. Koo, J. Yang}, Difference of weighted composition operators II, \textit{Integral Equ. Oper Theory}, \textbf{93} (2021), Paper No. 17, 19pp.

\bibitem{6} \textsc{B. Choe, T. Hosokawa, H. Koo}, Hilbert-Schmidt differences of composition operators on the Bergman space, \textit{Math. Z.}, \textbf{269} (2011), 751-775.

\bibitem{7} \textsc{B. Choe, H. Koo, I. Park}, Compact differences of composition operators on the Bergman spaces over the ball, \textit{Potential Anal.}, \textbf{40} (2014), 81-102.

\bibitem{8} \textsc{B. Choe, H. Koo, W. Smith}, Difference of composition operators over the half-plane, \textit{Trans. Amer. Math. Soc.}, \textbf{369} (2017), 3173-3205.

\bibitem{9} \textsc{C. Cowen, B. MacCluer}, \textit{Composition Operators on Spaces of Analytic Functions}, Studies in Advanced Mathematics, CRC Press, Boca Raton, FL, 1995.

\bibitem{10} \textsc{$\rm\check{Z}$. $\rm\check{C}$u$\rm\check{c}$kovi\'o, R. Zhao}, Weighted composition operators on the Bergman spaces, \textit{J. Lond. Math. Soc.}, \textbf{70}(2004), 499-511.

\bibitem{11} \textsc{$\rm\check{Z}$. $\rm\check{C}$u$\rm\check{c}$kovi\'o, R. Zhao}, Weighted composition operators between different weighted Bergman spaces and different Hardy spaces, \textit{Illinois J. Math.}, \textbf{51} (2007), 479-498.

\bibitem{12} \textsc{P. Duren}, \textit{Theory of $H^p$ Spaces}, Academic Press, New York-London, 1970, Reprint: Dover, Mineola, New York, 2000.

\bibitem{13} \textsc{F. Forelli}, The isometries of $H^p$. \textit{Canad. J. Math.}, \textbf{16} (1964), 721-728.

\bibitem{14} \textsc{C. Kolaski}, Isometries of weighted Bergman spaces, \textit{Canad. J. Math.}, \textbf{34} (1982), 910-915.

\bibitem{15} \textsc{H. Koo, M. Wang}, Joint Carleson measure and the difference of composition operators on $A_{\alpha}^p(\mathbb{B}^n)$, \textit{J. Math. Anal. Appl.}, \textbf{419} (2014), 1119-1142.

\bibitem{16} \textsc{C. Lo}, Norms and essential norms of differences of weighted composition operators, \textit{Mediterr. J. Math.}, \textbf{19} (2022), Paper No. 210, 15pp.

\bibitem{17} \textsc{D. Luecking}, Forward and reverse Carleson inequalities for functions in Bergman spaces and their derivatives, \textit{Amer. J. Math.}, \textbf{107} (1985), 85-111.

\bibitem{18} \textsc{D. Luecking}, Embedding theorems for spaces of analytic functions via Khinchine's inequality, \textit{Mich. Math. J.}, \textbf{40} (1993), 333-358.

\bibitem{19} \textsc{J. Moorhouse}, Compact differences of composition operators, \textit{J. Funct. Anal.}, \textbf{219} (2005), 70-92.

\bibitem{20} \textsc{I. Park}, Characterizations for the difference of weighted composition operators over the polydisk, \textit{J. Math. Anal. Appl.}, \textbf{533} (2024), 128011.

\bibitem{21} \textsc{E. Saukko}, Difference of composition operators between standard weighted Bergman spaces, \textit{J. Math. Anal. Appl.}, \textbf{381} (2011), 789-798.

\bibitem{22} \textsc{E. Saukko}, An application of atomic decomposition in Bergman spaces to the study of differences of composition operators, \textit{J. Funct. Anal.}, \textbf{262} (2012), 3872-3890.

\bibitem{23} \textsc{J. Shapiro}, \textit{Composition Operators and Classical Function Theory}, Springer-Verlag, New York, 1993.

\bibitem{24} \textsc{J. Shapiro, C. Sundberg}, Isolation amongst the composition operators, \textit{Pacific J. Math.}, \textbf{145} (1990), 117-152.

\bibitem{25} \textsc{Z. Wu}, A new characterization for Carleson measures and some applications, \textit{Integral Equ. Opera. Theory}, \textbf{71} (2011), 161-180.

\bibitem{26} \textsc{Z. Yang, Z. Zhou}, Atomic decomposition and composition operators on variable exponent Bergman spaces, \textit{Mediterr. J. Math.}, \textbf{19} (2022), Paper No. 10, 20pp.

\bibitem{27} \textsc{L. Zhang, Z. Zhou}, Hilbert-Schmidt differences of composition operators between the weighted Bergman spaces on the unit ball, \textit{Banach J. Math. Anal.}, \textit{7} (2013), 160-172.

\bibitem{28} \textsc{K. Zhu}, \textit{Spaces of Holomorphic Functions in the Unit Ball}, Springer-Verlag, New-York, 2005.

\bibitem{29} \textsc{K. Zhu}, \textit{Operator Theory in Function Spaces}, Second Edition. Mathematical Surveys and Monographs, American Mathematical Society, 2007.
	
	
	
\end{thebibliography}





%% If you have bib database file and want bibtex to generate the
%% bibitems, please use
%%
%%  \bibliographystyle{elsarticle-num-names} 
%%  \bibliography{<your bibdatabase>}

%% else use the following coding to input the bibitems directly in the
%% TeX file.

%% Refer following link for more details about bibliography and citations.
%% https://en.wikibooks.org/wiki/LaTeX/Bibliography_Management



%\endinput
\end{document}
%%
%% End of file `elsarticle-template-num-names.tex'.
